\documentclass[12pt,a4paper]{article}

% ── Packages ──────────────────────────────────────────────────────────────────
\usepackage[utf8]{inputenc}
\usepackage[romanian]{babel}
\usepackage{graphicx}
\usepackage{float}
\usepackage{listings}
\usepackage{xcolor}
\usepackage{geometry}
\usepackage{hyperref}
\usepackage{fancyhdr}
\usepackage{titlesec}
\usepackage{tocloft}
\usepackage{caption}
\usepackage{booktabs}
\usepackage{longtable}
\usepackage{array}
\usepackage{tikz}
\usepackage{pgf-umlcd}
\usepackage{todonotes}
\usepackage{mdframed}
\usepackage{enumitem}
\usepackage{multirow}
\usepackage{tabularx}
\usepackage{amsmath}
\usepackage{amssymb}

% ── Part formatting: no number, large bold centered, rule below, no page break ─
\titleclass{\part}{straight}
\titleformat{\part}[block]
  {\normalfont\Large\bfseries\centering}
  {}
  {0em}
  {}
  [\vspace{0.4ex}\rule{\textwidth}{0.8pt}]
\titlespacing*{\part}{0pt}{4ex}{2.5ex}

% ── Page layout ───────────────────────────────────────────────────────────────
\geometry{a4paper, left=2.5cm, right=2cm, top=2.5cm, bottom=2.5cm}

\pagestyle{fancy}
\fancyhf{}
\fancyhead[L]{GlobalBank DB — Baze de Date Distribuite}
\fancyhead[R]{MODBD IF 2025-2026}
\fancyfoot[C]{\thepage}

% ── SQL / Python listing style ─────────────────────────────────────────────────
\lstdefinestyle{sqlstyle}{
    language=SQL,
    basicstyle=\ttfamily\footnotesize,
    keywordstyle=\color{blue}\bfseries,
    commentstyle=\color{gray}\itshape,
    stringstyle=\color{red},
    showstringspaces=false,
    breaklines=true,
    frame=single,
    numbers=left,
    numberstyle=\tiny\color{gray},
    backgroundcolor=\color{gray!8},
    tabsize=2,
    captionpos=b
}
\lstdefinestyle{pystyle}{
    language=Python,
    basicstyle=\ttfamily\footnotesize,
    keywordstyle=\color{purple}\bfseries,
    commentstyle=\color{gray}\itshape,
    stringstyle=\color{teal},
    showstringspaces=false,
    breaklines=true,
    frame=single,
    numbers=left,
    numberstyle=\tiny\color{gray},
    backgroundcolor=\color{blue!4},
    tabsize=2,
    captionpos=b
}
\lstset{language={},basicstyle=\ttfamily\footnotesize,breaklines=true,frame=single,numbers=none,tabsize=2,captionpos=b}

% ── Screenshot placeholder ────────────────────────────────────────────────────
% Usage: \screenshot{subfolder/filename.png}{Caption}
\newcommand{\screenshot}[2]{%
  \begin{figure}[H]
    \centering
    \IfFileExists{../assets/screenshots/#1}{%
      \includegraphics[width=0.95\textwidth]{../assets/screenshots/#1}%
    }{%
      \begin{mdframed}[backgroundcolor=yellow!15, linecolor=orange, linewidth=1pt]
        \centering\vspace{0.5cm}
        \textcolor{orange}{\textbf{[SCREENSHOT NECESAR]}}\\[0.3cm]
        \texttt{#1}\\[0.3cm]
        \textit{#2}
        \vspace{0.5cm}
      \end{mdframed}%
    }%
    \caption{#2}
  \end{figure}%
}

% ── Diagram placeholder ───────────────────────────────────────────────────────
\newcommand{\diagramplaceholder}[2]{%
  \begin{figure}[H]
    \centering
    \IfFileExists{../assets/diagrams/#1}{%
      \includegraphics[width=0.95\textwidth]{../assets/diagrams/#1}%
    }{%
      \begin{mdframed}[backgroundcolor=blue!5, linecolor=blue!40, linewidth=1pt]
        \centering\vspace{1cm}
        \textcolor{blue!60}{\textbf{[DIAGRAMĂ]}}\\[0.3cm]
        \texttt{#1}\\[0.3cm]
        \textit{#2}
        \vspace{1cm}
      \end{mdframed}%
    }%
    \caption{#2}
  \end{figure}%
}

% ── Hyperref ──────────────────────────────────────────────────────────────────
\hypersetup{
    colorlinks=true, linkcolor=blue, urlcolor=cyan,
    pdftitle={GlobalBank DB — Proiect MODBD 2025-2026},
    pdfauthor={Sabău Eduard},
}

\graphicspath{{../assets/screenshots/}{../assets/diagrams/}}

% ── Title ─────────────────────────────────────────────────────────────────────
\title{
    \textbf{Metode de Optimizare și Distribuire în Baze de Date} \\[0.5em]
    \large Proiect: GlobalBank DB \\
    \normalsize Sistem Bancar Distribuit — Oracle 23ai Free (Docker)
}
\author{
    \textbf{Echipă}: Sabău Eduard, Liciu Ștefan \\
    \textbf{Grupa}: 506 \\
    \textbf{Coordonator}: Gabriela Mihai \\
    \textbf{Semestrul}: 2, 2025-2026
}

% ══════════════════════════════════════════════════════════════════════════════
\begin{document}
% ══════════════════════════════════════════════════════════════════════════════

\maketitle\thispagestyle{empty}

\begin{center}
\begin{tabular}{ll}
  \textbf{Baza de date}: & Oracle 23ai Free (\texttt{gvenzl/oracle-free:23-slim}) \\
  \textbf{Deployment}:   & Docker — 3 containere independente pe rețeaua \texttt{globalbank-net} \\
  \textbf{Porturi}:      & Central:1521 \quad București:1522 \quad Cluj:1523 \\
\end{tabular}
\end{center}

\newpage
\tableofcontents
\newpage

% ══════════════════════════════════════════════════════════════════════════════
\part{Modul Analiză (N1)}
% ══════════════════════════════════════════════════════════════════════════════

\section{Descrierea modelului și a obiectivelor aplicației}

\textbf{GlobalBank DB} este un sistem de baze de date distribuite destinat unei instituții bancare
cu sucursale în București și Cluj. Sistemul asigură:
\begin{itemize}
  \item Gestiunea distribuită a clienților, conturilor și tranzacțiilor bancare
  \item Transparența datelor prin vederi globale și triggere INSTEAD OF
  \item Optimizarea accesului local prin fragmentare orizontală și verticală
  \item Consistența cataloagelor prin replicare sincronă master→replică
\end{itemize}

Arhitectura rețelei de baze de date:
\begin{itemize}
  \item \textbf{Nod Central (\texttt{globalbank-central}:1521 — GLOBAL\_USER)}: coordonator, vederi globale, credite, profiluri clienți
  \item \textbf{Nod Local 1 (\texttt{globalbank-bucharest}:1522 — BUCHAREST\_USER)}: sucursala București, conturi B, tranzacții B
  \item \textbf{Nod Local 2 (\texttt{globalbank-cluj}:1523 — CLUJ\_USER)}: sucursala Cluj, conturi C, tranzacții C
\end{itemize}
Comunicarea inter-nod se realizează prin \textbf{DB Links} standard Oracle (\texttt{BUCHAREST\_LINK}, \texttt{CLUJ\_LINK}),
create cu \texttt{CREATE DATABASE LINK \ldots{} USING '(DESCRIPTION=\ldots{})'} (protocol TCP, rețea Docker bridge \texttt{globalbank-net}).

% ──────────────────────────────────────────────────────────────────────────────
\section{Diagramele bazei de date OLTP inițiale}

\subsection{Diagrama Entitate-Relație (ER)}

Baza de date OLTP inițială conține \textbf{10 entități independente}, cu normalizare FN3
și o relație many-to-many (\texttt{ANGAJAT\_CLIENT}).

\diagramplaceholder{er-diagram.png}{Diagrama Entitate-Relație — GlobalBank DB (10 entități, FN3)}

Entitățile și relațiile principale:

\begin{longtable}{lll}
\toprule
\textbf{Entitate} & \textbf{Atribute cheie} & \textbf{Relații} \\
\midrule
SUCURSALE & ID\_Sucursala (PK), Denumire, Oras & — \\
CLIENTI & ID\_Client, Nume, Prenume, CNP, Email, Scor\_Credit & — \\
CONTURI & ID\_Cont (PK), IBAN, Sold, Moneda & CLIENTI, TIPURI\_CONT, SUCURSALE \\
TRANZACTII & ID\_Tranzactie (PK), Data, Suma, Tip & CONTURI \\
CARDURI & ID\_Card (PK), Numar\_Card, Status & CONTURI \\
ANGAJATI & ID\_Angajat (PK), Nume, Functie, Salariu & SUCURSALE \\
ANGAJAT\_CLIENT & ID\_Angajat, ID\_Client, Rol & ANGAJATI $\times$ CLIENTI (\textbf{M:N}) \\
TIPURI\_CONT & ID\_Tip (PK), Denumire, Dobanda & — \\
CREDITE & ID\_Credit (PK), Suma\_Totala, Rata\_Lunara & CLIENTI \\
PLATI\_RATE & ID\_Plata (PK), Data\_Plata, Suma & CREDITE \\
\bottomrule
\end{longtable}

\textbf{Normalizare FN3}:
\begin{itemize}
  \item FN1: toate atributele sunt atomice
  \item FN2: toate atributele non-cheie depind de întreaga cheie primară
  \item FN3: nu există dependențe tranzitive (ex: Scor\_Credit depinde de ID\_Client, nu de altceva)
\end{itemize}

\subsection{Diagrama conceptuală OLTP}

\diagramplaceholder{conceptual-diagram.png}{Diagrama conceptuală a bazei de date OLTP inițiale}

% ──────────────────────────────────────────────────────────────────────────────
\section{Descrierea modului de distribuire}

Rețeaua de baze de date conține \textbf{3 containere Docker independente}, fiecare cu propria instanță Oracle 23ai Free:

\begin{center}
\begin{tabular}{lllll}
\toprule
\textbf{Container} & \textbf{Port} & \textbf{Schemă} & \textbf{Rol} & \textbf{PDB} \\
\midrule
\texttt{globalbank-central}   & 1521 & GLOBAL\_USER    & Nod Global / Central & \texttt{freepdb1} \\
\texttt{globalbank-bucharest} & 1522 & BUCHAREST\_USER & Nod Local 1          & \texttt{freepdb1} \\
\texttt{globalbank-cluj}      & 1523 & CLUJ\_USER      & Nod Local 2          & \texttt{freepdb1} \\
\bottomrule
\end{tabular}
\end{center}

Imaginea \texttt{gvenzl/oracle-free:23-slim} suportă nativ ARM64 (Apple Silicon) și x86-64.
Fiecare container are propriul proces Oracle (SID=FREE, PDB=freepdb1), volume de date separate
și este conectat doar prin DB Links standard pe rețeaua \texttt{globalbank-net}.

\textit{Notă}: O variantă alternativă de deployment folosește \textbf{Oracle ATP Free Tier} —
2 instanțe cloud (ATP1 pentru Central, ATP2 pentru nodurile locale în scheme separate).
În varianta ATP, DB Links se creează cu \texttt{DBMS\_CLOUD\_ADMIN} și folosesc protocol TCPS/mTLS.

% ──────────────────────────────────────────────────────────────────────────────
\section{Argumentarea deciziei de fragmentare}

\subsection{Fragmentare orizontală primară — relația CONTURI}

Schema bazei de date conține relația \texttt{CONTURI}, care stochează informațiile despre
conturile bancare gestionate de sistem.

\texttt{CONTURI}(\underline{ID\_Cont}, IBAN, ID\_Tip, Sold, Moneda, ID\_Sucursala, ID\_Client)

\begin{center}
\small
\begin{tabular}{c|l|c|c|c|c|c}
\toprule
ID\_Cont & IBAN & ID\_Tip & Sold & Moneda & ID\_Sucursala & ID\_Client \\
\midrule
1       & RO49GLOB0000001000000001 & 1 & 12000 & RON & 1 & 1 \\
2       & RO49GLOB0000001000000002 & 2 &  5500 & RON & 1 & 2 \\
3       & RO49GLOB0000001000000003 & 1 &  8200 & EUR & 1 & 3 \\
$\cdots$ & $\cdots$ & $\cdots$ & $\cdots$ & $\cdots$ & $\cdots$ & $\cdots$ \\
1000001 & RO49GLOB0001001000000001 & 1 &  9500 & RON & 2 & 1 \\
1000002 & RO49GLOB0001001000000002 & 3 & 15000 & RON & 2 & 2 \\
1000003 & RO49GLOB0001001000000003 & 2 &  4300 & EUR & 2 & 4 \\
\bottomrule
\end{tabular}
\end{center}

Atributul \texttt{ID\_Sucursala} a fost ales pentru fragmentare deoarece are un domeniu finit — $\{1, 2\}$
(1 = sucursala București, 2 = sucursala Cluj) — iar aplicațiile locale
accesează preponderent conturile sucursalei proprii.


\textbf{Legăturile owner--member}

Relația \texttt{CONTURI} conține cheia externă \texttt{ID\_Sucursala} și este \textit{owner}
în legăturile cu relațiile \textit{member} \texttt{TRANZACTII} și \texttt{CARDURI}, care vor
fi fragmentate orizontal derivat.

\begin{itemize}
  \item $owner(L_1) = \texttt{CONTURI}$, \quad $member(L_1) = \texttt{TRANZACTII}$
  \item $owner(L_2) = \texttt{CONTURI}$, \quad $member(L_2) = \texttt{CARDURI}$
\end{itemize}

\textbf{Aplicarea algoritmului de fragmentare orizontală primară (date ipotetice)}

\textit{Informații calitative - determinarea predicatelor simple}

Aplicațiile care accesează relația \texttt{CONTURI}:

Unde Q = {q1, q2, ..., qn} este mulțimea de aplicații utilizator
\begin{itemize}
  \item $q_1$: „Afișează conturile sucursalei București" $\rightarrow$ predicat $ID\_Sucursala = 1$
  \item $q_2$: „Afișează conturile sucursalei Cluj" $\rightarrow$ predicat $ID\_Sucursala = 2$
  \item $q_3$: „Procesează tranzacție nouă la București" $\rightarrow$ predicat $ID\_Sucursala = 1$
  \item $q_4$: „Procesează tranzacție nouă la Cluj" $\rightarrow$ predicat $ID\_Sucursala = 2$
  \item $q_5$: „Raport global: sold total pe toate conturile" $\rightarrow$ accesează toate tuplurile
\end{itemize}

Predicatele simple cele mai utilizate sunt:
\[
  Pr = \{p_1,\; p_2\}, \quad \text{unde} \quad
  p_1: ID\_Sucursala = 1, \qquad p_2: ID\_Sucursala = 2
\]

\textit{Informații cantitative}

Selectivitatea compusă — numărul de tupluri accesate de fiecare predicat:
\begin{itemize}
  \item $sel(ID\_Sucursala = 1) = 10$ (10 conturi la sucursala București)
  \item $sel(ID\_Sucursala = 2) = 10$ (10 conturi la sucursala Cluj)
\end{itemize}

Frecvența accesului:
\begin{itemize}
  \item $acc(q_1) = 40$ accesări/zi, \quad $acc(q_2) = 40$ accesări/zi
  \item $acc(q_3) = 30$ accesări/zi, \quad $acc(q_4) = 30$ accesări/zi
  \item $acc(q_5) = 5$ accesări/zi
\end{itemize}

\textit{Pasul 1 - Algoritmul COMP\_MIN}

Generează dintr-o mulțime Pr de predicate simple o mulțime Pr' completă și minimală.

Regula 1 - o relație sau un fragment este divizat in două sau mai multe subrelații, accesate diferit de cel puțin o aplicație.

\begin{description}
  \item[Input:] $R$: \texttt{CONTURI}; \quad $Pr = \{p_1, p_2\}$ unde $p_1: ID\_Sucursala=1$, $p_2: ID\_Sucursala=2$
  \item[Output:] $Pr'$: mulțime de predicate simple completă și minimală
    \begin{enumerate}
      \item Cautăm un predicat $p_i \in Pr$ care partiționează relația \texttt{CONTURI} conform regulii 1. Se alege $p_1: ID\_Sucursala=1$. Acesta partiționează \texttt{CONTURI} în:
        \begin{itemize}
          \item $f_1 = \sigma_{ID\_Sucursala=1}(\texttt{CONTURI})$ = \{conturi sucursala București\}
          \item $f_1' = \sigma_{ID\_Sucursala \neq 1}(\texttt{CONTURI})$ = \{conturi sucursala Cluj\}
        \end{itemize}
        $\Rightarrow$ $p_1$ este relevant și partiționează relația R. \\
        Inițializare:
        $Pr' \leftarrow \{p_1\}$, $Pr \leftarrow \{p_2\}$, $F \leftarrow \{f_1\}$
      \item (ciclu repeat-until): Se verifică $p_2: ID\_Sucursala=2$. Fragmentul $rest$ poate fi partiționat:
        \begin{itemize}
          \item $f_2 = \sigma_{ID\_Sucursala=2}(f_1')$ = \{conturi sucursala Cluj\}
          \item $f_1'' = \sigma_{ID\_Sucursala \neq 2}(f_1') = \emptyset$
        \end{itemize}
        Deoarece $domain(ID\_Sucursala) = \{1,2\}$, $p_2$ este relevant. \\
        Actualizare:
        $Pr' \leftarrow \{p_1, p_2\}$, $Pr \leftarrow \emptyset$, $F \leftarrow \{f_1, f_2\}$
        Verificam rezultatul obținut:
        \textbf{Verificare minimalitate}: Dacă eliminăm $p_1$, nu mai putem distinge conturile din București de cele din Cluj $\Rightarrow$ $p_1$ este relevant. Dacă eliminăm $p_2$, nu mai putem defini explicit fragmentul Cluj $\Rightarrow$ $p_2$ este relevant.
        \textbf{Verificare completitudine}: $Pr' = \{p_1, p_2\}$ este completă deoarece orice tuplu din \texttt{CONTURI} satisface $p_1$ sau $p_2$ (din \texttt{CHECK(ID\_Sucursala IN (1,2))}).
    \end{enumerate}
    \textbf{Rezultat}: $Pr' = \{p_1, p_2\}$ — mulțime completă și minimală.
\end{description}


\textit{Pasul 2 - Algoritmul PRIM\_ORIZ}

\begin{description}
  \item[Input:] $R$: \texttt{CONTURI}; \quad $Pr: \{p_1, p_2\}$
  \item[Output:] $M$: mulțime de predicate compuse ce determină fragmentele orizontale
    \begin{enumerate}
      \item $Pr' \leftarrow COMP\_MIN(\texttt{CONTURI}, Pr) = \{p_1, p_2\}$
      \item Se determină mulțimea $I$ de implicații. Deoarece $domain(ID\_Sucursala) = \{1, 2\}$:
        \begin{itemize}
          \item $i_1: (ID\_Sucursala=1) \Rightarrow \neg(ID\_Sucursala=2)$
          \item $i_2: \neg(ID\_Sucursala=1) \Rightarrow (ID\_Sucursala=2)$
        \end{itemize}
      \item Se determină mulțimea $M$ de predicate compuse. Din $Pr' = \{p_1, p_2\}$ se generează $2^2 = 4$ predicate compuse:
        \begin{itemize}
          \item $m_1: p_1 \land p_2 = (ID\_Sucursala=1) \land (ID\_Sucursala=2)$
          \item $m_2: p_1 \land \neg p_2 = (ID\_Sucursala=1) \land \neg(ID\_Sucursala=2)$
          \item $m_3: \neg p_1 \land p_2 = \neg(ID\_Sucursala=1) \land (ID\_Sucursala=2)$
          \item $m_4: \neg p_1 \land \neg p_2 = \neg(ID\_Sucursala=1) \land \neg(ID\_Sucursala=2)$
        \end{itemize}
      \item Se elimină predicatele compuse care nu au sens in raport $I$:
        \begin{itemize}
          \item $m_1$: $(ID\_Sucursala=1) \land (ID\_Sucursala=2)$ — \textbf{CONTRADICTIE} cu $i_1$
          \item $m_2$: $(ID\_Sucursala=1) \land \neg(ID\_Sucursala=2)$ — \textbf{VALID} (conform $i_1$, dacă $ID\_Sucursala=1$ atunci $\neg(ID\_Sucursala=2)$, condiție redundantă dar consistentă) $\Rightarrow$ \texttt{CONTURI\_B}
          \item $m_3$: $\neg(ID\_Sucursala=1) \land (ID\_Sucursala=2)$ — \textbf{VALID} (conform $i_2$, dacă $\neq(ID\_Sucursala=1)$ atunci $ID\_Sucursala=2$) $\Rightarrow$ \texttt{CONTURI\_C}
          \item $m_4$: $\neg(ID\_Sucursala=1) \land \neg(ID\_Sucursala=2)$ — \textbf{CONTRADICTORIU} cu $i_2$ $\Rightarrow$ eliminat
        \end{itemize}
    \end{enumerate}
    \textbf{Rezultat PRIM\_ORIZ}: $M = \{m_2, m_3\}$
\end{description}

\textbf{Obținerea fragmentelor orizontale primare}

Aplicând formula $R_i = \sigma_{F_i}(R)$, fragmentele orizontale ale relației \texttt{CONTURI} sunt:

\textit{Fragment 1: \texttt{CONTURI\_B} (stocat pe nodul Bucharest:1522)}
\[ CONTURI\_B = \sigma_{ID\_Sucursala=1}(CONTURI) \]

\begin{center}
\small
\begin{tabular}{c|l|c|c|c|c|c}
\toprule
ID\_Cont & IBAN & ID\_Tip & Sold & Moneda & ID\_Sucursala & ID\_Client \\
\midrule
1 & RO49GLOB0000001000000001 & 1 & 12000 & RON & 1 & 1 \\
2 & RO49GLOB0000001000000002 & 2 &  5500 & RON & 1 & 2 \\
3 & RO49GLOB0000001000000003 & 1 &  8200 & EUR & 1 & 3 \\
$\cdots$ & $\cdots$ & $\cdots$ & $\cdots$ & $\cdots$ & 1 & $\cdots$ \\
\bottomrule
\end{tabular}
\end{center}
(10 conturi total, $ID\_Cont \in [1, 999\,999]$)

\textit{Fragment 2: \texttt{CONTURI\_C} (stocat pe nodul Cluj:1523)}
\[ CONTURI\_C = \sigma_{ID\_Sucursala=2}(CONTURI) \]

\begin{center}
\small
\begin{tabular}{c|l|c|c|c|c|c}
\toprule
ID\_Cont & IBAN & ID\_Tip & Sold & Moneda & ID\_Sucursala & ID\_Client \\
\midrule
1000001 & RO49GLOB0001001000000001 & 1 &  9500 & RON & 2 & 1 \\
1000002 & RO49GLOB0001001000000002 & 3 & 15000 & RON & 2 & 2 \\
1000003 & RO49GLOB0001001000000003 & 2 &  4300 & EUR & 2 & 4 \\
$\cdots$ & $\cdots$ & $\cdots$ & $\cdots$ & $\cdots$ & 2 & $\cdots$ \\
\bottomrule
\end{tabular}
\end{center}
(10 conturi total, $ID\_Cont \in [1\,000\,001, 1\,999\,999]$)

\subsection{Fragmentare orizontală derivată — TRANZACTII și CARDURI}

Relațiile \texttt{TRANZACTII} și \texttt{CARDURI} sunt legate de \texttt{CONTURI} prin chei externe.
Este natural ca tranzacțiile și cardurile unui cont să fie stocate pe același nod cu contul respectiv,
eliminând join-urile cross-nod în operațiunile zilnice.

\subsubsection*{Relația TRANZACTII}

\texttt{TRANZACTII}(\underline{ID\_Tranzactie}, Data, Suma, Tip\_Tranzactie, ID\_Cont\_Sursa)

\begin{center}
\small
\begin{tabular}{c|c|c|l|c}
\toprule
ID\_Tranzactie & Data & Suma & Tip\_Tranzactie & ID\_Cont\_Sursa \\
\midrule
1       & 2024-01-15 &  500 & Plata     & 1 \\
2       & 2024-01-16 & 1200 & Depunere  & 2 \\
3       & 2024-01-17 &  300 & Retragere & 3 \\
$\cdots$ & $\cdots$ & $\cdots$ & $\cdots$ & $\cdots$ \\
1000001 & 2024-01-15 &  800 & Plata     & 1000001 \\
1000002 & 2024-01-16 & 2000 & Depunere  & 1000002 \\
\bottomrule
\end{tabular}
\end{center}

\textbf{Obținerea fragmentelor orizontale derivate — TRANZACTII}

Fie $L_1$ legătura dintre \texttt{CONTURI} și \texttt{TRANZACTII} prin atributul \texttt{ID\_Cont\_Sursa}:
\[
  owner(L_1) = \texttt{CONTURI}, \qquad member(L_1) = \texttt{TRANZACTII}
\]

\[
  TRANZACTII_i = TRANZACTII \ltimes CONTURI_i, \quad i \in \{B, C\}
\]

\textit{Fragment 1: \texttt{TRANZACTII\_B} (stocat pe Bucharest:1522)}
\[ TRANZACTII\_B = TRANZACTII \ltimes CONTURI\_B \]
Deoarece $CONTURI\_B$ conține conturile cu $ID\_Cont \leq 999\,999$, semijoin-ul selectează
$ID\_Cont\_Sursa \leq 999\,999$:

\begin{center}
\small
\begin{tabular}{c|c|c|l|c}
\toprule
ID\_Tranzactie & Data & Suma & Tip\_Tranzactie & ID\_Cont\_Sursa \\
\midrule
1 & 2024-01-15 &  500 & Plata     & 1 \\
2 & 2024-01-16 & 1200 & Depunere  & 2 \\
3 & 2024-01-17 &  300 & Retragere & 3 \\
$\cdots$ & $\cdots$ & $\cdots$ & $\cdots$ & $\cdots$ \\
\bottomrule
\end{tabular}
\end{center}
(20 tranzacții, $ID\_Cont\_Sursa \in [1, 999\,999]$)

\textit{Fragment 2: \texttt{TRANZACTII\_C} (stocat pe Cluj:1523)}
\[ TRANZACTII\_C = TRANZACTII \ltimes CONTURI\_C \]

\begin{center}
\small
\begin{tabular}{c|c|c|l|c}
\toprule
ID\_Tranzactie & Data & Suma & Tip\_Tranzactie & ID\_Cont\_Sursa \\
\midrule
1000001 & 2024-01-15 &  800 & Plata     & 1000001 \\
1000002 & 2024-01-16 & 2000 & Depunere  & 1000002 \\
1000003 & 2024-01-17 &  450 & Retragere & 1000003 \\
$\cdots$ & $\cdots$ & $\cdots$ & $\cdots$ & $\cdots$ \\
\bottomrule
\end{tabular}
\end{center}
(20 tranzacții, $ID\_Cont\_Sursa \in [1\,000\,001, 1\,999\,999]$)

\subsubsection*{Relația CARDURI}

\texttt{CARDURI}(\underline{ID\_Card}, Numar\_Card, Data\_Expirare, Status, ID\_Cont)

Fie $L_2$ legătura dintre \texttt{CONTURI} și \texttt{CARDURI} prin atributul \texttt{ID\_Cont}:
\[
  owner(L_2) = \texttt{CONTURI}, \qquad member(L_2) = \texttt{CARDURI}
\]

Fragmentele derivate:
\[
  CARDURI\_B = CARDURI \ltimes CONTURI\_B \quad (ID\_Cont \in [1, 999\,999])
\]
\[
  CARDURI\_C = CARDURI \ltimes CONTURI\_C \quad (ID\_Cont \in [1\,000\,001, 1\,999\,999])
\]

\subsection{Fragmentare verticală — relația CLIENTI}

Se considera relația: \texttt{CLIENTI} cu următoarele atribute:

\texttt{CLIENTI}(\underline{ID\_Client}, Nume, Prenume, CNP, Email, Telefon, Scor\_Credit)

\begin{center}
\small
\begin{tabular}{c|l|l|l|l|l|c}
\toprule
ID\_Client & Nume & Prenume & CNP & Email & Telefon & Scor\_Credit \\
\midrule
1 & Popescu & Ion   & 1900101400001 & ion.pop@mail.ro    & 0721000001 & 720 \\
2 & Ionescu & Maria & 2900215400002 & maria.ion@mail.ro  & 0731000002 & 650 \\
3 & Popa    & Gh.   & 1850303400003 & gh.popa@mail.ro    & 0741000003 & 580 \\
$\cdots$ & $\cdots$ & $\cdots$ & $\cdots$ & $\cdots$ & $\cdots$ & $\cdots$ \\
\bottomrule
\end{tabular}
\end{center}

Notăm atributele:

\begin{center}
\begin{tabular}{l|l}
\toprule
Atribut & Descriere \\
\midrule
$A_1 = ID\_Client$ & Cheia primară (inclusă în ambele fragmente) \\
$A_2 = Nume$       & Numele clientului \\
$A_3 = Prenume$    & Prenumele clientului \\
$A_4 = CNP$        & Cod Numeric Personal \\
$A_5 = Email$      & Adresa de email \\
$A_6 = Telefon$    & Numărul de telefon \\
$A_7 = Scor\_Credit$ & Scorul de credit \\
\bottomrule
\end{tabular}
\end{center}

\textbf{Aplicarea algoritmului de fragmentare verticală (date ipotetice)}

\textit{Următoarele cereri accesează relația \texttt{CLIENTI}:}

\begin{itemize}
  \item $Q_1$ — Verificare identitate client la sucursală; atribute accesate: $A_1, A_2, A_3, A_4$
\begin{lstlisting}[language=SQL]
SELECT id_client, nume, prenume, cnp FROM clienti WHERE id_client = :id;
\end{lstlisting}
  \item $Q_2$ — Calculare/verificare scor de credit (aplicație centrală); atribute: $A_1, A_5, A_6, A_7$
\begin{lstlisting}[language=SQL]
SELECT id_client, email, telefon, scor_credit FROM clienti WHERE id_client = :id;
\end{lstlisting}
  \item $Q_3$ — Raport global cu toate informațiile clientului; atribute: $A_1, A_2, A_3, A_4, A_5, A_6, A_7$
\begin{lstlisting}[language=SQL]
SELECT * FROM clienti;
\end{lstlisting}
  \item $Q_4$ — Căutare client după date de contact; atribute: $A_1, A_2, A_3, A_5, A_6$
\begin{lstlisting}[language=SQL]
SELECT id_client, nume, prenume, email, telefon FROM clienti WHERE email = :email;
\end{lstlisting}
\end{itemize}

\textit{Matricea VA}

Funcția $use(q_i, A_j)$: 1 dacă $A_j$ este referit de $q_i$ , 0 altfel.

\begin{center}
\begin{tabular}{l|cccccc}
\toprule
 & $A_2$ (Nume) & $A_3$ (Prenume) & $A_4$ (CNP) & $A_5$ (Email) & $A_6$ (Telefon) & $A_7$ (Scor\_Credit) \\
\midrule
$Q_1$ & 1 & 1 & 1 & 0 & 0 & 0 \\
$Q_2$ & 0 & 0 & 0 & 1 & 1 & 1 \\
$Q_3$ & 1 & 1 & 1 & 1 & 1 & 1 \\
$Q_4$ & 1 & 1 & 0 & 1 & 1 & 0 \\
\bottomrule
\end{tabular}
\end{center}

\textit{Frecvențele de acces ale cererilor (pe toate stațiile):}

\begin{center}
\begin{tabular}{l|c|l}
\toprule
Cerere & $acc(q_k)$ (accesări/zi) & Justificare \\
\midrule
$Q_1$ & 40 & Verificare frecventă la sucursale locale \\
$Q_2$ & 20 & Scoring centralizat \\
$Q_3$ &  5 & Raportare globală \\
$Q_4$ & 15 & Căutare client după date de contact \\
\bottomrule
\end{tabular}
\end{center}

\textbf{Matricea afinității atributelor}

Afinitatea dintre atributele $A_i$ și $A_j$ se calculează după formula:
\[
  aff(A_i, A_j) = \sum_{\substack{k \in K \\ \forall S_l}} ref_l(q_k) \cdot acc_l(q_k),
  \quad \text{unde } K = \{k \mid use(q_k, A_i) = 1 \land use(q_k, A_j) = 1\}
\]

\begin{itemize}
  \item $aff(A_2, A_3)$: $Q_1(40) + Q_3(5) + Q_4(15) = \mathbf{60}$
  \item $aff(A_2, A_4)$: $Q_1(40) + Q_3(5) = \mathbf{45}$
  \item $aff(A_2, A_5)$: $Q_3(5) + Q_4(15) = \mathbf{20}$
  \item $aff(A_2, A_6)$: $Q_3(5) + Q_4(15) = \mathbf{20}$
  \item $aff(A_2, A_7)$: $Q_3(5) = \mathbf{5}$
  \item $aff(A_3, A_4)$: $Q_1(40) + Q_3(5) = \mathbf{45}$
  \item $aff(A_3, A_5)$: $Q_3(5) + Q_4(15) = \mathbf{20}$ \quad $aff(A_3, A_6) = \mathbf{20}$ \quad $aff(A_3, A_7) = \mathbf{5}$
  \item $aff(A_4, A_5)$: $Q_3(5) = \mathbf{5}$ \quad $aff(A_4, A_6) = \mathbf{5}$ \quad $aff(A_4, A_7) = \mathbf{5}$
  \item $aff(A_5, A_6)$: $Q_2(20) + Q_3(5) + Q_4(15) = \mathbf{40}$
  \item $aff(A_5, A_7)$: $Q_2(20) + Q_3(5) = \mathbf{25}$ \quad $aff(A_6, A_7) = \mathbf{25}$
\end{itemize}

Valorile diagonalei principale: $aff(A_2,A_2)=60$, $aff(A_3,A_3)=60$, $aff(A_4,A_4)=45$,
$aff(A_5,A_5)=40$, $aff(A_6,A_6)=40$, $aff(A_7,A_7)=25$.

\textbf{Matricea de afinitate $AA$}:

\begin{center}
\begin{tabular}{l|rrrrrr}
\toprule
$AA$ & $A_2$ & $A_3$ & $A_4$ & $A_5$ & $A_6$ & $A_7$ \\
\midrule
$A_2$ & 60 & 60 & 45 & 20 & 20 &  5 \\
$A_3$ & 60 & 60 & 45 & 20 & 20 &  5 \\
$A_4$ & 45 & 45 & 45 &  5 &  5 &  5 \\
$A_5$ & 20 & 20 &  5 & 40 & 40 & 25 \\
$A_6$ & 20 & 20 &  5 & 40 & 40 & 25 \\
$A_7$ &  5 &  5 &  5 & 25 & 25 & 25 \\
\bottomrule
\end{tabular}
\end{center}

% ──────────────────────────────────────────────────────────────────────────────
\section{Verificarea corectitudinii fragmentărilor}

\subsection{Fragmentare orizontală primară}

\begin{enumerate}
  \item \textbf{Completitudinea}: Este dată de algoritmul fragmentării ce determină o mulțime completă și minimală de predicate. Orice tuplu din relația \texttt{CONTURI} se regăsește în cel puțin unul dintre fragmente. Deoarece $domain(ID\_Sucursala) = \{1, 2\}$ și constrângerea \texttt{CHECK} este impusă prin DDL, orice rând satisface $p_1$ sau $p_2$:
  \[ CONTURI\_B \cup CONTURI\_C = CONTURI \]
  \item \textbf{Reconstrucția}: Relația globală \texttt{CONTURI} poate fi obținută din fragmentele sale prin operatorul de reuniune (implementat ca \texttt{UNION ALL} în vederea \texttt{CONTURI\_GLOBAL}):
  \[ CONTURI = CONTURI\_B \cup CONTURI\_C \]
  \item \textbf{Disjuncția}: Fragmentele sunt disjuncte deoarece predicatele compuse $m_2$ și $m_3$ se exclud reciproc — un cont nu poate aparține simultan sucursalei 1 și sucursalei 2:
  \[ CONTURI\_B \cap CONTURI\_C = \emptyset \]
\end{enumerate}

\subsection{Fragmentare orizontala derivata}

\begin{enumerate}
  \item \textbf{Completitudinea}: Pentru fiecare tuplu $t$ al relatiei \texttt{TRANZACTII} , exista un tuplu $t'$ al relatiei \texttt{CONTURI} corespondent. Orice tranzacție are $ID\_Cont\_Sursa$ cheie externă validă; contul corespunzător se află în $CONTURI\_B$ sau $CONTURI\_C$, deci:
  \[ TRANZACTII\_B \cup TRANZACTII\_C = TRANZACTII \]
  \item \textbf{Reconstrucția}: Relația globală \texttt{TRANZACTII} poate fi reconstituită prin reuniunea fragmentelor (implementată ca \texttt{UNION ALL} în vederea \texttt{TRANZACTII\_GLOBAL}):
  \[ TRANZACTII = TRANZACTII\_B \cup TRANZACTII\_C \]
  \item \textbf{Disjuncția}: Graful relațiilor și legăturilor dintre cele doua relatii este simplu (o singură legătură owner--member între \texttt{CONTURI} și \texttt{TRANZACTII}). Domeniile $ID\_Cont\_Sursa$ sunt disjuncte ($[1, 999\,999]$ față de $[1\,000\,001, 1\,999\,999]$):
  \[ TRANZACTII\_B \cap TRANZACTII\_C = \emptyset \]
\end{enumerate}

Proprietățile de completitudine, reconstrucție și disjuncție sunt identice pentru fragmentele derivate ale relației \texttt{CARDURI}.


\subsection{Fragmentare verticala}

\begin{enumerate}
  \item \textbf{Completitudinea}: Fiecare atribut al relației se găseste intr-un singur fragment în urma partiționării acesteia :
  \[ Attr(CLIENTI\_ID) \cup Attr(CLIENTI\_PROFIL) = Attr(CLIENTI) \]
  \item \textbf{Reconstrucția}: Se realizează prin operatorul de compunere după cheia primară($ID\_Client$), implementat în vederea \texttt{CLIENTI\_GLOBAL}:
  \[ CLIENTI = CLIENTI\_ID \bowtie_{ID\_Client} CLIENTI\_PROFIL \]
  \item \textbf{Disjuncția}: Cheia primară este replicată in fiecare fragment, așadar se poate considera ca fragmentele satisfac proprietatea de disjunctie:
  \[ Attr(CLIENTI\_ID) \cap Attr(CLIENTI\_PROFIL) = \{ID\_Client\} \]
\end{enumerate}
% ──────────────────────────────────────────────────────────────────────────────
\section{Argumentarea deciziei de replicare}

\subsection{Relații replicate}

Replicarea este justificată pentru relațiile de tip \textit{catalog} cu volum mic,
accesate frecvent de pe nodurile locale și referite prin FK de fragmentele orizontale.

\begin{center}
{\small
\begin{tabular}{p{2.6cm}p{2.8cm}p{5.4cm}p{2.6cm}}
\toprule
\textbf{Relație} & \textbf{Tip replicare} & \textbf{Motivație} & \textbf{Noduri destinație} \\
\midrule
TIPURI\_CONT & Sincronă, single-master & Catalog de 4 rânduri (Curent, Economii, Depozit, Credit);
  citit la fiecare interogare de cont;
  FK local din CONTURI\_B/C necesită copie locală pentru a evita latența DB Link &
  București, Cluj \\
\bottomrule
\end{tabular}
}
\end{center}

\textbf{Mecanism}: triggerul \texttt{TRG\_REPLICATE\_TIPURI\_CONT} propagă orice INSERT/UPDATE/DELETE
de pe nodul Central (master) simultan pe ambele noduri locale prin DB Links, în aceeași tranzacție.
Nodurile locale sunt \textit{read-only} pentru această relație — orice modificare trebuie efectuată
exclusiv pe Central.

\subsection{Relații stocate pe o singură stație}

Relațiile care nu sunt nici fragmentate, nici replicate rămân integral pe un singur nod,
justificat prin natura datelor sau prin apartenența lor exclusivă la un centru de responsabilitate.

\begin{center}
{\small
\begin{tabular}{p{3.8cm}p{2cm}p{7cm}}
\toprule
\textbf{Relație} & \textbf{Stație} & \textbf{Motivație} \\
\midrule
SUCURSALE & Central & Tabel de referință cu 2 înregistrări (București ID=1, Cluj ID=2);
  accesat rar și exclusiv de GLOBAL\_USER; volum prea mic pentru a justifica fragmentarea,
  iar replicarea pe noduri locale ar crea dependențe inutile. \\
CREDITE & Central & Creditele sunt aprobate și gestionate centralizat; corelarea directă cu
  CLIENTI\_PROFIL (același nod) elimină necesitatea DB Links pentru constrângeri de integritate. \\
PLATI\_RATE & Central & Ratele sunt înregistrate centralizat; FK direct la CREDITE (aceeași stație)
  permite constrângere DDL fără triggere cross-nod. \\
LOG\_ACCES & Central & Jurnal de audit global; inserările vin din toate nodurile via DB Link;
  centralizarea simplifică conformitatea și raportarea de securitate. \\
ANGAJATI\_B & București & Angajați specifici sucursalei București; relevanți exclusiv
  pentru operațiunile locale; transferul pe Central ar adăuga trafic de rețea inutil. \\
ANGAJAT\_CLIENT\_B & București & Relație M:N locală între angajații și clienții gestionați
  de sucursala București. \\
ANGAJATI\_C & Cluj & Angajați specifici sucursalei Cluj; structură simetrică cu
  ANGAJATI\_B, pe nodul corespunzător. \\
ANGAJAT\_CLIENT\_C & Cluj & Relație M:N locală Cluj, simetrică cu ANGAJAT\_CLIENT\_B. \\
\bottomrule
\end{tabular}
}
\end{center}

% ──────────────────────────────────────────────────────────────────────────────
\section{Schemele conceptuale locale}

\subsection{Schema locală București (BUCHAREST\_USER — container Bucharest:1522)}

\begin{center}
\begin{tabular}{ll}
\toprule
\textbf{Tabel} & \textbf{Descriere} \\
\midrule
TIPURI\_CONT & Copie replica a catalogului (replicat de pe Central) \\
CLIENTI\_ID & Fragment vertical: date de identificare (Nume, CNP) \\
CONTURI\_B & Fragment orizontal: conturi cu ID\_Sucursala=1 \\
TRANZACTII\_B & Fragment orizontal derivat: tranzacții pentru CONTURI\_B \\
CARDURI\_B & Date locale: carduri atașate conturilor B \\
ANGAJATI\_B & Date locale: angajați sucursala București \\
ANGAJAT\_CLIENT\_B & Relație M:N locală: angajat gestionează client \\
\bottomrule
\end{tabular}
\end{center}

Secvențe: \texttt{SEQ\_CLIENT\_B} (1–999\,999), \texttt{SEQ\_CONT\_B} (1–999\,999), \texttt{SEQ\_TRANZ\_B} (1–999\,999)

\screenshot{bucharest-local-schema.png}{Schemă locala București}


\subsection{Schema locală Cluj (CLUJ\_USER — container Cluj:1523)}

\begin{center}
\begin{tabular}{ll}
\toprule
\textbf{Tabel} & \textbf{Descriere} \\
\midrule
TIPURI\_CONT & Copie replica a catalogului (replicat de pe Central) \\
CLIENTI\_ID & Fragment vertical: date de identificare (Nume, CNP) \\
CONTURI\_C & Fragment orizontal: conturi cu ID\_Sucursala=1 \\
TRANZACTII\_C & Fragment orizontal derivat: tranzacții pentru CONTURI\_C \\
CARDURI\_C & Date locale: carduri atașate conturilor C \\
ANGAJATI\_C & Date locale: angajați sucursala Cluj \\
ANGAJAT\_CLIENT\_C & Relație M:N locală: angajat gestionează client \\
\bottomrule
\end{tabular}
\end{center}

Secvențe: \texttt{SEQ\_CLIENT\_C} \texttt{SEQ\_CONT\_C}  \texttt{SEQ\_TRANZ\_C} , toate porning de la 1\,000\,001 pentru unicitate globală

\screenshot{cluj-local-schema.png}{Schemă locala Cluj}

% ──────────────────────────────────────────────────────────────────────────────
\section{Lista constrângerilor de integritate}

\subsection{Unicitate}

\textbf{Unicitate locală} (per fragment):

\begin{center}
\begin{tabular}{llll}
\toprule
\textbf{Relație} & \textbf{Coloană} & \textbf{Tip} & \textbf{Nod} \\
\midrule
CONTURI\_B     & IBAN      & UNIQUE & București \\
CONTURI\_C     & IBAN      & UNIQUE & Cluj \\
CLIENTI\_ID    & CNP       & UNIQUE & București, Cluj \\
TIPURI\_CONT   & Denumire  & UNIQUE & Central, București, Cluj \\
SUCURSALE      & Denumire  & UNIQUE & Central \\
CARDURI\_B     & Numar\_Card & UNIQUE & București \\
CARDURI\_C     & Numar\_Card & UNIQUE & Cluj \\
CLIENTI\_ID    & Email     & UNIQUE & București, Cluj \\
\bottomrule
\end{tabular}
\end{center}

\textbf{Unicitate globală pe fragmente orizontale}:
Garantată prin ranguri de secvențe non-suprapuse:

\begin{center}
{\footnotesize
\begin{tabular}{llll}
\toprule
\textbf{Secvență} & \textbf{Interval} & \textbf{Tabel} & \textbf{Nod} \\
\midrule
SEQ\_CONT\_B        & [1,\ 999\,999]              & CONTURI\_B    & București \\
SEQ\_CONT\_C        & [1\,000\,001, 1\,999\,999]  & CONTURI\_C    & Cluj \\
SEQ\_TRANZ\_B       & [1,\ 999\,999]              & TRANZACTII\_B & București \\
SEQ\_TRANZ\_C       & [1\,000\,001, 1\,999\,999]  & TRANZACTII\_C & Cluj \\
SEQ\_CLIENT\_GLOBAL & [2\,000\,001, $\infty$)     & CLIENTI\_ID   & ambele noduri \\
\bottomrule
\end{tabular}
}
\end{center}

Fără coordonare centralizată, fiecare fragment produce ID-uri unice global.
Suplimentar, triggerul \texttt{trg\_iof\_conturi\_global} verifică unicitatea IBAN
cross-fragment (CONTURI\_B $\cup$ CONTURI\_C) înainte de orice INSERT.

\textbf{Unicitate globală pentru fragmente verticale}:
\texttt{CLIENTI\_ID.CNP} este UNIQUE pe fiecare nod local $\Rightarrow$ unicitatea globală a CNP
este garantată de triggerul INSTEAD OF care inserează simultan pe ambii noduri
(dacă CNP există deja pe oricare nod, INSERT eșuează prin constrângerea locală).

\subsection{Cheie primară}

\begin{center}
\begin{tabular}{lll}
\toprule
\textbf{Tabel} & \textbf{Cheie primară} & \textbf{Nivel} \\
\midrule
CONTURI\_B / CONTURI\_C & ID\_Cont & Local \\
TRANZACTII\_B / C & ID\_Tranzactie & Local \\
CLIENTI\_ID & ID\_Client & Local \\
CLIENTI\_PROFIL & ID\_Client & Central \\
TIPURI\_CONT & ID\_Tip & Local + Central \\
CREDITE & ID\_Credit & Central \\
PLATI\_RATE & ID\_Plata & Central \\
SUCURSALE & ID\_Sucursala & Central \\
\bottomrule
\end{tabular}
\end{center}

\subsection{Cheie externă}

\textbf{Locale (DDL, în același container)}:

\begin{center}
{\small
\begin{tabular}{p{4.2cm}p{2.5cm}p{3cm}p{2.5cm}p{1.5cm}}
\toprule
\textbf{Tabel sursă} & \textbf{Coloană} & \textbf{Tabel referit} & \textbf{Coloană} & \textbf{Nod} \\
\midrule
CONTURI\_B       & ID\_Tip        & TIPURI\_CONT  & ID\_Tip        & București \\
CONTURI\_B       & ID\_Client     & CLIENTI\_ID   & ID\_Client     & București \\
CONTURI\_C       & ID\_Tip        & TIPURI\_CONT  & ID\_Tip        & Cluj \\
CONTURI\_C       & ID\_Client     & CLIENTI\_ID   & ID\_Client     & Cluj \\
CARDURI\_B       & ID\_Cont       & CONTURI\_B    & ID\_Cont       & București \\
CARDURI\_C       & ID\_Cont       & CONTURI\_C    & ID\_Cont       & Cluj \\
TRANZACTII\_B    & ID\_Cont\_Sursa & CONTURI\_B   & ID\_Cont       & București \\
TRANZACTII\_C    & ID\_Cont\_Sursa & CONTURI\_C   & ID\_Cont       & Cluj \\
ANGAJAT\_CLIENT\_B & ID\_Angajat  & ANGAJATI\_B   & ID\_Angajat    & București \\
ANGAJAT\_CLIENT\_C & ID\_Angajat  & ANGAJATI\_C   & ID\_Angajat    & Cluj \\
PLATI\_RATE      & ID\_Credit     & CREDITE       & ID\_Credit     & Central \\
\bottomrule
\end{tabular}
}
\end{center}

\textbf{Cross-nod (simulate prin triggere — Oracle nu suportă DDL FK inter-instance)}:

\begin{center}
{\small
\begin{tabular}{p{4.5cm}p{3.2cm}p{5.5cm}p{1.5cm}}
\toprule
\textbf{Trigger} & \textbf{Eveniment} & \textbf{Constrângere simulată} & \textbf{Nod} \\
\midrule
trg\_bef\_ins\_credite &
  BEFORE INSERT ON CREDITE &
  CREDITE.ID\_Client $\in$ CLIENTI\_ID @BUCHAREST\_LINK &
  Central \\
trg\_iof\_clienti\_global &
  INSTEAD OF INSERT ON CLIENTI\_GLOBAL &
  Inserare simultană pe ambele noduri; eșec dacă CNP există pe oricare nod &
  Central \\
trg\_iof\_conturi\_global &
  INSTEAD OF INSERT ON CONTURI\_GLOBAL &
  Verificare IBAN global unic în CONTURI\_B $\cup$ CONTURI\_C înainte de INSERT &
  Central \\
\bottomrule
\end{tabular}
}
\end{center}

\subsection{Validare (CHECK constraints)}

\begin{center}
{\small
\begin{tabular}{p{3cm}p{7cm}p{4cm}}
\toprule
\textbf{Tabel} & \textbf{Constrângere CHECK} & \textbf{Scop} \\
\midrule
CONTURI\_B      & \texttt{CHECK(ID\_Sucursala = 1)}                              & Puritate fragment orizontal \\
CONTURI\_C      & \texttt{CHECK(ID\_Sucursala = 2)}                              & Puritate fragment orizontal \\
TRANZACTII\_B/C & \texttt{CHECK(Tip\_Tranzactie IN ('DEBIT','CREDIT','TRANSFER'))} & Validare domeniu \\
TIPURI\_CONT    & \texttt{CHECK(Dobanda BETWEEN 0 AND 100)}                       & Validare domeniu \\
CLIENTI\_PROFIL & \texttt{CHECK(Scor\_Credit BETWEEN 0 AND 999)}                  & Validare scoring \\
\bottomrule
\end{tabular}
}
\end{center}

% ══════════════════════════════════════════════════════════════════════════════
\setcounter{section}{0}
\part{Modul Implementare Baze de Date (N2)}
% ══════════════════════════════════════════════════════════════════════════════

\section{Crearea bazelor de date și utilizatorilor}

\subsection{GLOBAL\_USER pe Central (Docker)}
\begin{lstlisting}
CREATE USER GLOBAL_USER IDENTIFIED BY "SecurePass123!";

GRANT CONNECT, RESOURCE      TO GLOBAL_USER;
GRANT CREATE TABLE           TO GLOBAL_USER;
GRANT CREATE VIEW            TO GLOBAL_USER;
GRANT CREATE SEQUENCE        TO GLOBAL_USER;
GRANT CREATE TRIGGER         TO GLOBAL_USER;
GRANT CREATE PROCEDURE       TO GLOBAL_USER;
GRANT CREATE DATABASE LINK   TO GLOBAL_USER;
GRANT UNLIMITED TABLESPACE   TO GLOBAL_USER;

-- Allow GLOBAL_USER to create synonyms and query data dictionary
GRANT CREATE SYNONYM         TO GLOBAL_USER;
GRANT SELECT ANY DICTIONARY  TO GLOBAL_USER;

COMMIT;

\end{lstlisting}
\screenshot{create-global-user.png}{GLOBAL\_USER creat cu succes pe containerul Central}

\subsection{BUCHAREST\_USER pe Bucharest (Docker)}
\begin{lstlisting}
-- Create schema user
CREATE USER BUCHAREST_USER IDENTIFIED BY "SecurePass123!";

GRANT CONNECT, RESOURCE    TO BUCHAREST_USER;
GRANT CREATE TABLE         TO BUCHAREST_USER;
GRANT CREATE VIEW          TO BUCHAREST_USER;
GRANT CREATE SEQUENCE      TO BUCHAREST_USER;
GRANT CREATE TRIGGER       TO BUCHAREST_USER;
GRANT CREATE PROCEDURE     TO BUCHAREST_USER;
GRANT UNLIMITED TABLESPACE TO BUCHAREST_USER;

COMMIT;

\end{lstlisting}
\screenshot{create-bucharest-user.png}{BUCHAREST\_USER creat cu succes pe containerul Bucharest}

\subsection{CLUJ\_USER pe Cluj (Docker)}
\begin{lstlisting}
    -- Create schema user
CREATE USER CLUJ_USER IDENTIFIED BY "SecurePass123!";

GRANT CONNECT, RESOURCE    TO CLUJ_USER;
GRANT CREATE TABLE         TO CLUJ_USER;
GRANT CREATE VIEW          TO CLUJ_USER;
GRANT CREATE SEQUENCE      TO CLUJ_USER;
GRANT CREATE TRIGGER       TO CLUJ_USER;
GRANT CREATE PROCEDURE     TO CLUJ_USER;
GRANT UNLIMITED TABLESPACE TO CLUJ_USER;

COMMIT;
\end{lstlisting}
\screenshot{create-cluj-user.png}{CLUJ\_USER creat cu succes pe containerul Cluj}

\subsection{DB Links (GLOBAL\_USER pe Central)}
\begin{lstlisting}
-- BUCHAREST_LINK
CREATE DATABASE LINK BUCHAREST_LINK
  CONNECT TO BUCHAREST_USER IDENTIFIED BY "SecurePass123!"
  USING '(DESCRIPTION=(ADDRESS=(PROTOCOL=TCP)(HOST=bucharest)(PORT=1521))(CONNECT_DATA=(SERVICE_NAME=freepdb1)))';

-- CLUJ_LINK
CREATE DATABASE LINK CLUJ_LINK
  CONNECT TO CLUJ_USER IDENTIFIED BY "SecurePass123!"
  USING '(DESCRIPTION=(ADDRESS=(PROTOCOL=TCP)(HOST=cluj)(PORT=1521))(CONNECT_DATA=(SERVICE_NAME=freepdb1)))';

-- Verify both links
SELECT 'BUCHAREST_LINK OK - remote user: ' || USER AS status FROM DUAL@BUCHAREST_LINK;
SELECT 'CLUJ_LINK OK - remote user: '      || USER AS status FROM DUAL@CLUJ_LINK;
\end{lstlisting}
\screenshot{db-links.png}{BUCHAREST\_LINK și CLUJ\_LINK verificate cu SELECT USER FROM DUAL@...}

% ──────────────────────────────────────────────────────────────────────────────
\section{Crearea relațiilor și fragmentelor}

\subsection{Schema centrală (GLOBAL\_USER — container Central:1521)}
\begin{lstlisting}

-- SUCURSALE
CREATE TABLE SUCURSALE (
    ID_Sucursala  NUMBER(4)     PRIMARY KEY,
    Nume          VARCHAR2(100) NOT NULL,
    Oras          VARCHAR2(50)  NOT NULL,
    Adresa        VARCHAR2(200) NOT NULL,
    CONSTRAINT chk_sucursala_oras CHECK (Oras IN ('Bucuresti', 'Cluj'))
);

INSERT INTO SUCURSALE VALUES (1, 'GlobalBank Bucuresti Centru', 'Bucuresti', 'Bulevardul Unirii 10, Sector 4');
INSERT INTO SUCURSALE VALUES (2, 'GlobalBank Cluj Napoca',      'Cluj',      'Strada Memorandumului 5, Cluj-Napoca');

-- TIPURI_CONT master si replicata pe nodurile Bucuresti/Cluj prin triggeri
CREATE TABLE TIPURI_CONT (
    ID_Tip    NUMBER(4)    PRIMARY KEY,
    Denumire  VARCHAR2(50) NOT NULL UNIQUE,
    Dobanda   NUMBER(5,2)  NOT NULL,
    CONSTRAINT chk_dobanda CHECK (Dobanda >= 0 AND Dobanda <= 100)
);

INSERT INTO TIPURI_CONT VALUES (1, 'Curent',   0.00);
INSERT INTO TIPURI_CONT VALUES (2, 'Economii', 3.50);
INSERT INTO TIPURI_CONT VALUES (3, 'Depozit',  5.00);
INSERT INTO TIPURI_CONT VALUES (4, 'Credit',  15.00);

-- CLIENTI_PROFIL
-- SEQ_CLIENT_GLOBAL folosita de un trigger INSTEAD OF pe view-ul CLIENTI_GLOBAL
CREATE SEQUENCE SEQ_CLIENT_GLOBAL START WITH 2000001 MAXVALUE 9999999 NOCACHE NOCYCLE;

CREATE TABLE CLIENTI_PROFIL (
    ID_Client    NUMBER(10)    PRIMARY KEY,
    Email        VARCHAR2(150) UNIQUE,
    Telefon      VARCHAR2(15),
    Scor_Credit  NUMBER(3)     DEFAULT 500,
    CONSTRAINT chk_scor  CHECK (Scor_Credit BETWEEN 0 AND 999),
    CONSTRAINT chk_email CHECK (Email LIKE '%@%.%')
);

-- CREDITE
CREATE SEQUENCE SEQ_CREDIT START WITH 2000001 MAXVALUE 9999999 NOCACHE NOCYCLE;

CREATE TABLE CREDITE (
    ID_Credit    NUMBER(10)   PRIMARY KEY,
    Suma_Totala  NUMBER(15,2) NOT NULL,
    Rata_Lunara  NUMBER(10,2) NOT NULL,
    ID_Client    NUMBER(10)   NOT NULL,
    CONSTRAINT chk_credit_suma CHECK (Suma_Totala > 0),
    CONSTRAINT chk_credit_rata CHECK (Rata_Lunara > 0)
);

-- PLATI_RATE
CREATE SEQUENCE SEQ_PLATA START WITH 2000001 MAXVALUE 9999999 NOCACHE NOCYCLE;

CREATE TABLE PLATI_RATE (
    ID_Plata    NUMBER(10)   PRIMARY KEY,
    Data_Plata  DATE         DEFAULT SYSDATE NOT NULL,
    Suma        NUMBER(10,2) NOT NULL,
    ID_Credit   NUMBER(10)   NOT NULL,
    CONSTRAINT chk_plata_suma  CHECK (Suma > 0),
    CONSTRAINT fk_plata_credit FOREIGN KEY (ID_Credit) REFERENCES CREDITE(ID_Credit)
);

-- LOG_ACCES
CREATE SEQUENCE SEQ_LOG START WITH 1 MAXVALUE 99999999 NOCACHE NOCYCLE;

CREATE TABLE LOG_ACCES (
    ID_Log     NUMBER(10)    PRIMARY KEY,
    Utilizator VARCHAR2(50)  NOT NULL,
    Data_Ora   TIMESTAMP     DEFAULT SYSTIMESTAMP NOT NULL,
    Actiune    VARCHAR2(200) NOT NULL
);

COMMIT;

\end{lstlisting}
\screenshot{central-schema.png}{Tabele centrale create: SUCURSALE, TIPURI\_CONT, CLIENTI\_PROFIL, CREDITE, PLATI\_RATE, LOG\_ACCES}

\subsection{Fragment orizontal — București (BUCHAREST\_USER — container Bucharest:1522)}
\begin{lstlisting}
-- TIPURI_CONT replica
CREATE TABLE TIPURI_CONT (
    ID_Tip    NUMBER(4)    PRIMARY KEY,
    Denumire  VARCHAR2(50) NOT NULL UNIQUE,
    Dobanda   NUMBER(5,2)  NOT NULL,
    CONSTRAINT chk_dobanda_b CHECK (Dobanda >= 0 AND Dobanda <= 100)
);

-- CLIENTI_ID
CREATE SEQUENCE SEQ_CLIENT_B START WITH 1 MAXVALUE 999999 NOCACHE NOCYCLE;

CREATE TABLE CLIENTI_ID (
    ID_Client  NUMBER(10)    PRIMARY KEY,
    Nume       VARCHAR2(100) NOT NULL,
    Prenume    VARCHAR2(100) NOT NULL,
    CNP        CHAR(13)      NOT NULL UNIQUE,
    CONSTRAINT chk_cnp_len_b CHECK (LENGTH(CNP) = 13)
);

-- CONTURI_B
CREATE SEQUENCE SEQ_CONT_B START WITH 1 MAXVALUE 999999 NOCACHE NOCYCLE;

CREATE TABLE CONTURI_B (
    ID_Cont      NUMBER(10)   PRIMARY KEY,
    IBAN         VARCHAR2(34) NOT NULL UNIQUE,
    ID_Tip       NUMBER(4)    NOT NULL,
    Sold         NUMBER(15,2) DEFAULT 0,
    Moneda       VARCHAR2(3)  DEFAULT 'RON',
    ID_Sucursala NUMBER(4)    DEFAULT 1 NOT NULL,
    ID_Client    NUMBER(10)   NOT NULL,
    CONSTRAINT chk_conturi_b_suc CHECK (ID_Sucursala = 1),
    CONSTRAINT chk_moneda_b      CHECK (Moneda IN ('RON','EUR','USD')),
    CONSTRAINT chk_sold_b        CHECK (Sold >= 0),
    CONSTRAINT fk_conturi_b_cli  FOREIGN KEY (ID_Client)  REFERENCES CLIENTI_ID(ID_Client),
    CONSTRAINT fk_conturi_b_tip  FOREIGN KEY (ID_Tip)     REFERENCES TIPURI_CONT(ID_Tip)
);

-- TRANZACTII_B
CREATE SEQUENCE SEQ_TRANZ_B START WITH 1 MAXVALUE 999999 NOCACHE NOCYCLE;

CREATE TABLE TRANZACTII_B (
    ID_Tranzactie  NUMBER(10)   PRIMARY KEY,
    Data           DATE         DEFAULT SYSDATE NOT NULL,
    Suma           NUMBER(15,2) NOT NULL,
    Tip_Tranzactie VARCHAR2(20) NOT NULL,
    ID_Cont_Sursa  NUMBER(10)   NOT NULL,
    CONSTRAINT chk_tranz_b_suma CHECK (Suma > 0),
    CONSTRAINT chk_tranz_b_tip  CHECK (Tip_Tranzactie IN ('DEBIT','CREDIT','TRANSFER')),
    CONSTRAINT fk_tranz_b_cont  FOREIGN KEY (ID_Cont_Sursa) REFERENCES CONTURI_B(ID_Cont)
);

-- T018: CARDURI_B
CREATE TABLE CARDURI_B (
    ID_Card        NUMBER(10)   PRIMARY KEY,
    Numar_Card     CHAR(16)     NOT NULL UNIQUE,
    Data_Expirare  DATE         NOT NULL,
    Status         VARCHAR2(10) DEFAULT 'ACTIV',
    ID_Cont        NUMBER(10)   NOT NULL,
    CONSTRAINT chk_card_b_status CHECK (Status IN ('ACTIV','BLOCAT','EXPIRAT')),
    CONSTRAINT fk_card_b_cont    FOREIGN KEY (ID_Cont) REFERENCES CONTURI_B(ID_Cont)
);

-- T019: ANGAJATI_B
CREATE TABLE ANGAJATI_B (
    ID_Angajat   NUMBER(10)    PRIMARY KEY,
    Nume         VARCHAR2(100) NOT NULL,
    Functie      VARCHAR2(50)  NOT NULL,
    Salariu      NUMBER(10,2)  NOT NULL,
    ID_Sucursala NUMBER(4)     DEFAULT 1 NOT NULL,
    CONSTRAINT chk_ang_b_suc CHECK (ID_Sucursala = 1),
    CONSTRAINT chk_ang_b_sal CHECK (Salariu > 0)
);

-- ANGAJAT_CLIENT_B (M:N junction)
CREATE TABLE ANGAJAT_CLIENT_B (
    ID_Angajat    NUMBER(10)   NOT NULL,
    ID_Client     NUMBER(10)   NOT NULL,
    Rol           VARCHAR2(30) DEFAULT 'GESTIONAR',
    Data_Asignare DATE         DEFAULT SYSDATE,
    CONSTRAINT pk_ac_b     PRIMARY KEY (ID_Angajat, ID_Client),
    CONSTRAINT fk_ac_b_ang FOREIGN KEY (ID_Angajat) REFERENCES ANGAJATI_B(ID_Angajat),
    CONSTRAINT fk_ac_b_cli FOREIGN KEY (ID_Client)  REFERENCES CLIENTI_ID(ID_Client)
);

COMMIT;
\end{lstlisting}
\screenshot{bucharest-schema.png}{Tabele București create: CONTURI\_B, TRANZACTII\_B, CLIENTI\_ID și altele}

\subsection{Fragment orizontal — Cluj (CLUJ\_USER — container Cluj:1523)}
\begin{lstlisting}
-- TIPURI_CONT replica
CREATE TABLE TIPURI_CONT (
    ID_Tip    NUMBER(4)    PRIMARY KEY,
    Denumire  VARCHAR2(50) NOT NULL UNIQUE,
    Dobanda   NUMBER(5,2)  NOT NULL,
    CONSTRAINT chk_dobanda_c CHECK (Dobanda >= 0 AND Dobanda <= 100)
);

-- CLIENTI_ID
-- Sequence range 1000000-1999999
CREATE SEQUENCE SEQ_CLIENT_C START WITH 1000001 MAXVALUE 1999999 NOCACHE NOCYCLE;

CREATE TABLE CLIENTI_ID (
    ID_Client  NUMBER(10)    PRIMARY KEY,
    Nume       VARCHAR2(100) NOT NULL,
    Prenume    VARCHAR2(100) NOT NULL,
    CNP        CHAR(13)      NOT NULL UNIQUE,
    CONSTRAINT chk_cnp_len_c CHECK (LENGTH(CNP) = 13)
);

-- CONTURI_C
CREATE SEQUENCE SEQ_CONT_C START WITH 1000001 MAXVALUE 1999999 NOCACHE NOCYCLE;

CREATE TABLE CONTURI_C (
    ID_Cont      NUMBER(10)   PRIMARY KEY,
    IBAN         VARCHAR2(34) NOT NULL UNIQUE,
    ID_Tip       NUMBER(4)    NOT NULL,
    Sold         NUMBER(15,2) DEFAULT 0,
    Moneda       VARCHAR2(3)  DEFAULT 'RON',
    ID_Sucursala NUMBER(4)    DEFAULT 2 NOT NULL,
    ID_Client    NUMBER(10)   NOT NULL,
    CONSTRAINT chk_conturi_c_suc CHECK (ID_Sucursala = 2),
    CONSTRAINT chk_moneda_c      CHECK (Moneda IN ('RON','EUR','USD')),
    CONSTRAINT chk_sold_c        CHECK (Sold >= 0),
    CONSTRAINT fk_conturi_c_cli  FOREIGN KEY (ID_Client)  REFERENCES CLIENTI_ID(ID_Client),
    CONSTRAINT fk_conturi_c_tip  FOREIGN KEY (ID_Tip)     REFERENCES TIPURI_CONT(ID_Tip)
);

-- TRANZACTII_C
CREATE SEQUENCE SEQ_TRANZ_C START WITH 1000001 MAXVALUE 1999999 NOCACHE NOCYCLE;

CREATE TABLE TRANZACTII_C (
    ID_Tranzactie  NUMBER(10)   PRIMARY KEY,
    Data           DATE         DEFAULT SYSDATE NOT NULL,
    Suma           NUMBER(15,2) NOT NULL,
    Tip_Tranzactie VARCHAR2(20) NOT NULL,
    ID_Cont_Sursa  NUMBER(10)   NOT NULL,
    CONSTRAINT chk_tranz_c_suma CHECK (Suma > 0),
    CONSTRAINT chk_tranz_c_tip  CHECK (Tip_Tranzactie IN ('DEBIT','CREDIT','TRANSFER')),
    CONSTRAINT fk_tranz_c_cont  FOREIGN KEY (ID_Cont_Sursa) REFERENCES CONTURI_C(ID_Cont)
);

-- CARDURI_C
CREATE TABLE CARDURI_C (
    ID_Card        NUMBER(10)   PRIMARY KEY,
    Numar_Card     CHAR(16)     NOT NULL UNIQUE,
    Data_Expirare  DATE         NOT NULL,
    Status         VARCHAR2(10) DEFAULT 'ACTIV',
    ID_Cont        NUMBER(10)   NOT NULL,
    CONSTRAINT chk_card_c_status CHECK (Status IN ('ACTIV','BLOCAT','EXPIRAT')),
    CONSTRAINT fk_card_c_cont    FOREIGN KEY (ID_Cont) REFERENCES CONTURI_C(ID_Cont)
);

-- ANGAJATI_C
CREATE TABLE ANGAJATI_C (
    ID_Angajat   NUMBER(10)    PRIMARY KEY,
    Nume         VARCHAR2(100) NOT NULL,
    Functie      VARCHAR2(50)  NOT NULL,
    Salariu      NUMBER(10,2)  NOT NULL,
    ID_Sucursala NUMBER(4)     DEFAULT 2 NOT NULL,
    CONSTRAINT chk_ang_c_suc CHECK (ID_Sucursala = 2),
    CONSTRAINT chk_ang_c_sal CHECK (Salariu > 0)
);

-- ANGAJAT_CLIENT_C (M:N)
CREATE TABLE ANGAJAT_CLIENT_C (
    ID_Angajat    NUMBER(10)   NOT NULL,
    ID_Client     NUMBER(10)   NOT NULL,
    Rol           VARCHAR2(30) DEFAULT 'GESTIONAR',
    Data_Asignare DATE         DEFAULT SYSDATE,
    CONSTRAINT pk_ac_c     PRIMARY KEY (ID_Angajat, ID_Client),
    CONSTRAINT fk_ac_c_ang FOREIGN KEY (ID_Angajat) REFERENCES ANGAJATI_C(ID_Angajat),
    CONSTRAINT fk_ac_c_cli FOREIGN KEY (ID_Client)  REFERENCES CLIENTI_ID(ID_Client)
);

COMMIT;

\end{lstlisting}
\screenshot{cluj-schema.png}{Tabele Cluj create: CONTURI\_C, TRANZACTII\_C, CLIENTI\_ID și altele}

% ──────────────────────────────────────────────────────────────────────────────
\section{Popularea cu date}

\begin{lstlisting}
-- CLIENTI_PROFIL -- 10 clients (IDs 1-10)
INSERT INTO CLIENTI_PROFIL VALUES (1,  'popescu.ion@gmail.com',         '0721234567', 750);
INSERT INTO CLIENTI_PROFIL VALUES (2,  'ionescu.maria@gmail.com',       '0732345678', 820);
INSERT INTO CLIENTI_PROFIL VALUES (3,  'georgescu.mihai@yahoo.ro',      '0743456789', 680);
INSERT INTO CLIENTI_PROFIL VALUES (4,  'stan.elena@gmail.com',          '0754567890', 900);
INSERT INTO CLIENTI_PROFIL VALUES (5,  'dumitrescu.andrei@hotmail.com', '0765678901', 450);
INSERT INTO CLIENTI_PROFIL VALUES (6,  'constantin.ana@gmail.com',      '0776789012', 600);
INSERT INTO CLIENTI_PROFIL VALUES (7,  'popa.dan@yahoo.ro',             '0787890123', 720);
INSERT INTO CLIENTI_PROFIL VALUES (8,  'moldovan.ioana@gmail.com',      '0798901234', 850);
INSERT INTO CLIENTI_PROFIL VALUES (9,  'radu.cristian@yahoo.ro',        '0709012345', 550);
INSERT INTO CLIENTI_PROFIL VALUES (10, 'barbu.simona@gmail.com',        '0710123456', 780);

-- CREDITE
INSERT INTO CREDITE VALUES (SEQ_CREDIT.NEXTVAL, 25000.00,  550.00, 1);
INSERT INTO CREDITE VALUES (SEQ_CREDIT.NEXTVAL, 50000.00, 1200.00, 2);
INSERT INTO CREDITE VALUES (SEQ_CREDIT.NEXTVAL, 15000.00,  350.00, 3);
INSERT INTO CREDITE VALUES (SEQ_CREDIT.NEXTVAL, 80000.00, 1800.00, 4);
INSERT INTO CREDITE VALUES (SEQ_CREDIT.NEXTVAL, 10000.00,  250.00, 5);

-- PLATI_RATE
INSERT INTO PLATI_RATE VALUES (SEQ_PLATA.NEXTVAL, SYSDATE - 60,  550.00, 2000001);
INSERT INTO PLATI_RATE VALUES (SEQ_PLATA.NEXTVAL, SYSDATE - 30,  550.00, 2000001);
INSERT INTO PLATI_RATE VALUES (SEQ_PLATA.NEXTVAL, SYSDATE,        550.00, 2000001);
INSERT INTO PLATI_RATE VALUES (SEQ_PLATA.NEXTVAL, SYSDATE - 60, 1200.00, 2000002);
INSERT INTO PLATI_RATE VALUES (SEQ_PLATA.NEXTVAL, SYSDATE - 30, 1200.00, 2000002);
INSERT INTO PLATI_RATE VALUES (SEQ_PLATA.NEXTVAL, SYSDATE,      1200.00, 2000002);
INSERT INTO PLATI_RATE VALUES (SEQ_PLATA.NEXTVAL, SYSDATE - 60,  350.00, 2000003);
INSERT INTO PLATI_RATE VALUES (SEQ_PLATA.NEXTVAL, SYSDATE - 30,  350.00, 2000003);
INSERT INTO PLATI_RATE VALUES (SEQ_PLATA.NEXTVAL, SYSDATE,        350.00, 2000003);
INSERT INTO PLATI_RATE VALUES (SEQ_PLATA.NEXTVAL, SYSDATE - 60, 1800.00, 2000004);
INSERT INTO PLATI_RATE VALUES (SEQ_PLATA.NEXTVAL, SYSDATE - 30, 1800.00, 2000004);
INSERT INTO PLATI_RATE VALUES (SEQ_PLATA.NEXTVAL, SYSDATE,      1800.00, 2000004);
INSERT INTO PLATI_RATE VALUES (SEQ_PLATA.NEXTVAL, SYSDATE - 60,  250.00, 2000005);
INSERT INTO PLATI_RATE VALUES (SEQ_PLATA.NEXTVAL, SYSDATE - 30,  250.00, 2000005);
INSERT INTO PLATI_RATE VALUES (SEQ_PLATA.NEXTVAL, SYSDATE,        250.00, 2000005);

-- LOG_ACCES
INSERT INTO LOG_ACCES VALUES (SEQ_LOG.NEXTVAL, 'GLOBAL_USER',    SYSTIMESTAMP - INTERVAL '5'  DAY, 'Conectare la nodul central');
INSERT INTO LOG_ACCES VALUES (SEQ_LOG.NEXTVAL, 'BUCHAREST_USER', SYSTIMESTAMP - INTERVAL '4'  DAY, 'INSERT in CONTURI_B');
INSERT INTO LOG_ACCES VALUES (SEQ_LOG.NEXTVAL, 'CLUJ_USER',      SYSTIMESTAMP - INTERVAL '3'  DAY, 'INSERT in CONTURI_C');
INSERT INTO LOG_ACCES VALUES (SEQ_LOG.NEXTVAL, 'GLOBAL_USER',    SYSTIMESTAMP - INTERVAL '1'  DAY, 'SELECT din CONTURI_GLOBAL');
INSERT INTO LOG_ACCES VALUES (SEQ_LOG.NEXTVAL, 'ADMIN',          SYSTIMESTAMP - INTERVAL '12' HOUR, 'Verificare DB Links');

COMMIT;
\end{lstlisting}
\screenshot{populate-central.png}{10 CLIENTI\_PROFIL, 5 CREDITE, 15 PLATI\_RATE inserate pe Central}

\begin{lstlisting}
-- TIPURI_CONT
INSERT INTO TIPURI_CONT VALUES (1, 'Curent',   0.00);
INSERT INTO TIPURI_CONT VALUES (2, 'Economii', 3.50);
INSERT INTO TIPURI_CONT VALUES (3, 'Depozit',  5.00);
INSERT INTO TIPURI_CONT VALUES (4, 'Credit',  15.00);

-- CLIENTI_ID
INSERT INTO CLIENTI_ID VALUES (1,  'Popescu',      'Ion',     '1850101400123');
INSERT INTO CLIENTI_ID VALUES (2,  'Ionescu',      'Maria',   '2900215350456');
INSERT INTO CLIENTI_ID VALUES (3,  'Georgescu',    'Mihai',   '1751205180789');
INSERT INTO CLIENTI_ID VALUES (4,  'Stan',         'Elena',   '2680430120012');
INSERT INTO CLIENTI_ID VALUES (5,  'Dumitrescu',   'Andrei',  '1920714230345');
INSERT INTO CLIENTI_ID VALUES (6,  'Constantin',   'Ana',     '2870922560678');
INSERT INTO CLIENTI_ID VALUES (7,  'Popa',         'Dan',     '1780305090901');
INSERT INTO CLIENTI_ID VALUES (8,  'Moldovan',     'Ioana',   '2950816410234');
INSERT INTO CLIENTI_ID VALUES (9,  'Radu',         'Cristian','1631128270567');
INSERT INTO CLIENTI_ID VALUES (10, 'Barbu',        'Simona',  '2710309190890');

-- CONTURI_B
INSERT INTO CONTURI_B VALUES (1,  'RO49GLOB0000000000000001', 1, 5000.00,  'RON', 1, 1);
INSERT INTO CONTURI_B VALUES (2,  'RO49GLOB0000000000000002', 2, 12500.00, 'RON', 1, 2);
INSERT INTO CONTURI_B VALUES (3,  'RO49GLOB0000000000000003', 1,   800.00, 'RON', 1, 3);
INSERT INTO CONTURI_B VALUES (4,  'RO49GLOB0000000000000004', 3, 50000.00, 'RON', 1, 4);
INSERT INTO CONTURI_B VALUES (5,  'RO49GLOB0000000000000005', 1,  2300.00, 'RON', 1, 5);
INSERT INTO CONTURI_B VALUES (6,  'RO49GLOB0000000000000006', 2,  8900.00, 'EUR', 1, 6);
INSERT INTO CONTURI_B VALUES (7,  'RO49GLOB0000000000000007', 4,     0.00, 'RON', 1, 7);
INSERT INTO CONTURI_B VALUES (8,  'RO49GLOB0000000000000008', 1, 15600.00, 'RON', 1, 8);
INSERT INTO CONTURI_B VALUES (9,  'RO49GLOB0000000000000009', 3, 75000.00, 'RON', 1, 9);
INSERT INTO CONTURI_B VALUES (10, 'RO49GLOB0000000000000010', 2,  4200.00, 'EUR', 1, 10);

-- TRANZACTII_B
INSERT INTO TRANZACTII_B VALUES (1,  SYSDATE - 30, 1000.00, 'CREDIT',   1);
INSERT INTO TRANZACTII_B VALUES (2,  SYSDATE - 15,  500.00, 'DEBIT',    1);
INSERT INTO TRANZACTII_B VALUES (3,  SYSDATE - 25, 2000.00, 'CREDIT',   2);
INSERT INTO TRANZACTII_B VALUES (4,  SYSDATE - 10, 1500.00, 'DEBIT',    2);
INSERT INTO TRANZACTII_B VALUES (5,  SYSDATE - 20,  300.00, 'CREDIT',   3);
INSERT INTO TRANZACTII_B VALUES (6,  SYSDATE -  5,  200.00, 'DEBIT',    3);
INSERT INTO TRANZACTII_B VALUES (7,  SYSDATE - 22, 5000.00, 'CREDIT',   4);
INSERT INTO TRANZACTII_B VALUES (8,  SYSDATE -  8, 2000.00, 'DEBIT',    4);
INSERT INTO TRANZACTII_B VALUES (9,  SYSDATE - 18,  800.00, 'CREDIT',   5);
INSERT INTO TRANZACTII_B VALUES (10, SYSDATE -  3,  600.00, 'TRANSFER', 5);
INSERT INTO TRANZACTII_B VALUES (11, SYSDATE - 28, 1500.00, 'CREDIT',   6);
INSERT INTO TRANZACTII_B VALUES (12, SYSDATE - 12, 1000.00, 'DEBIT',    6);
INSERT INTO TRANZACTII_B VALUES (13, SYSDATE - 15, 2500.00, 'CREDIT',   7);
INSERT INTO TRANZACTII_B VALUES (14, SYSDATE -  6, 1200.00, 'DEBIT',    7);
INSERT INTO TRANZACTII_B VALUES (15, SYSDATE - 20, 3500.00, 'CREDIT',   8);
INSERT INTO TRANZACTII_B VALUES (16, SYSDATE -  9, 1800.00, 'DEBIT',    8);
INSERT INTO TRANZACTII_B VALUES (17, SYSDATE - 24, 8000.00, 'CREDIT',   9);
INSERT INTO TRANZACTII_B VALUES (18, SYSDATE - 11, 3000.00, 'DEBIT',    9);
INSERT INTO TRANZACTII_B VALUES (19, SYSDATE - 16,  900.00, 'CREDIT',   10);
INSERT INTO TRANZACTII_B VALUES (20, SYSDATE -  4,  700.00, 'TRANSFER', 10);

-- CARDURI_B
INSERT INTO CARDURI_B VALUES (1, '4532015112830366', DATE '2028-12-31', 'ACTIV',  1);
INSERT INTO CARDURI_B VALUES (2, '4916338506082832', DATE '2027-06-30', 'ACTIV',  2);
INSERT INTO CARDURI_B VALUES (3, '4539578763621486', DATE '2029-03-31', 'ACTIV',  3);
INSERT INTO CARDURI_B VALUES (4, '4485275742308325', DATE '2026-09-30', 'BLOCAT', 4);
INSERT INTO CARDURI_B VALUES (5, '4716114061269756', DATE '2028-05-31', 'ACTIV',  5);

-- ANGAJATI_B
INSERT INTO ANGAJATI_B VALUES (1, 'Stoica Bogdan',   'Director',           8500.00, 1);
INSERT INTO ANGAJATI_B VALUES (2, 'Neagu Alina',     'Consultant Bancar',  4500.00, 1);
INSERT INTO ANGAJATI_B VALUES (3, 'Preda Marius',    'Casier',             3200.00, 1);
INSERT INTO ANGAJATI_B VALUES (4, 'Dima Cristina',   'Analist Credite',    5200.00, 1);
INSERT INTO ANGAJATI_B VALUES (5, 'Manea Florin',    'Responsabil Carduri',4800.00, 1);

-- ANGAJAT_CLIENT_B
INSERT INTO ANGAJAT_CLIENT_B VALUES (1, 1,  'DIRECTOR',  SYSDATE - 180);
INSERT INTO ANGAJAT_CLIENT_B VALUES (2, 3,  'GESTIONAR', SYSDATE - 120);
INSERT INTO ANGAJAT_CLIENT_B VALUES (2, 4,  'GESTIONAR', SYSDATE - 90);
INSERT INTO ANGAJAT_CLIENT_B VALUES (3, 5,  'GESTIONAR', SYSDATE - 60);
INSERT INTO ANGAJAT_CLIENT_B VALUES (4, 2,  'ANALIST',   SYSDATE - 150);

COMMIT;

ALTER SEQUENCE SEQ_CLIENT_B INCREMENT BY 10; SELECT SEQ_CLIENT_B.NEXTVAL FROM DUAL; ALTER SEQUENCE SEQ_CLIENT_B INCREMENT BY 1;
ALTER SEQUENCE SEQ_CONT_B   INCREMENT BY 10; SELECT SEQ_CONT_B.NEXTVAL   FROM DUAL; ALTER SEQUENCE SEQ_CONT_B   INCREMENT BY 1;
ALTER SEQUENCE SEQ_TRANZ_B  INCREMENT BY 20; SELECT SEQ_TRANZ_B.NEXTVAL  FROM DUAL; ALTER SEQUENCE SEQ_TRANZ_B  INCREMENT BY 1;

\end{lstlisting}
\screenshot{populate-bucharest.png}{59 rânduri inserate pe nodul București}

\begin{lstlisting}
-- TIPURI_CONT
INSERT INTO TIPURI_CONT VALUES (1, 'Curent',   0.00);
INSERT INTO TIPURI_CONT VALUES (2, 'Economii', 3.50);
INSERT INTO TIPURI_CONT VALUES (3, 'Depozit',  5.00);
INSERT INTO TIPURI_CONT VALUES (4, 'Credit',  15.00);

-- CLIENTI_ID
INSERT INTO CLIENTI_ID VALUES (1,  'Popescu',    'Ion',      '1850101400123');
INSERT INTO CLIENTI_ID VALUES (2,  'Ionescu',    'Maria',    '2900215350456');
INSERT INTO CLIENTI_ID VALUES (3,  'Georgescu',  'Mihai',    '1751205180789');
INSERT INTO CLIENTI_ID VALUES (4,  'Stan',       'Elena',    '2680430120012');
INSERT INTO CLIENTI_ID VALUES (5,  'Dumitrescu', 'Andrei',   '1920714230345');
INSERT INTO CLIENTI_ID VALUES (6,  'Constantin', 'Ana',      '2870922560678');
INSERT INTO CLIENTI_ID VALUES (7,  'Popa',       'Dan',      '1780305090901');
INSERT INTO CLIENTI_ID VALUES (8,  'Moldovan',   'Ioana',    '2950816410234');
INSERT INTO CLIENTI_ID VALUES (9,  'Radu',       'Cristian', '1631128270567');
INSERT INTO CLIENTI_ID VALUES (10, 'Barbu',      'Simona',   '2710309190890');

-- CONTURI_C
INSERT INTO CONTURI_C VALUES (1000001, 'RO49GLOB0000001000000001', 1,  3200.00, 'RON', 2, 1);
INSERT INTO CONTURI_C VALUES (1000002, 'RO49GLOB0000001000000002', 2,  9800.00, 'RON', 2, 2);
INSERT INTO CONTURI_C VALUES (1000003, 'RO49GLOB0000001000000003', 1,  1500.00, 'RON', 2, 3);
INSERT INTO CONTURI_C VALUES (1000004, 'RO49GLOB0000001000000004', 3, 35000.00, 'RON', 2, 4);
INSERT INTO CONTURI_C VALUES (1000005, 'RO49GLOB0000001000000005', 2,  7600.00, 'EUR', 2, 5);
INSERT INTO CONTURI_C VALUES (1000006, 'RO49GLOB0000001000000006', 1,  4100.00, 'RON', 2, 6);
INSERT INTO CONTURI_C VALUES (1000007, 'RO49GLOB0000001000000007', 2, 22000.00, 'RON', 2, 7);
INSERT INTO CONTURI_C VALUES (1000008, 'RO49GLOB0000001000000008', 4,     0.00, 'RON', 2, 8);
INSERT INTO CONTURI_C VALUES (1000009, 'RO49GLOB0000001000000009', 1,   890.00, 'RON', 2, 9);
INSERT INTO CONTURI_C VALUES (1000010, 'RO49GLOB0000001000000010', 3, 45000.00, 'RON', 2, 10);

-- TRANZACTII_C
INSERT INTO TRANZACTII_C VALUES (1000001, SYSDATE - 29, 1200.00, 'CREDIT',   1000001);
INSERT INTO TRANZACTII_C VALUES (1000002, SYSDATE - 14,  600.00, 'DEBIT',    1000001);
INSERT INTO TRANZACTII_C VALUES (1000003, SYSDATE - 26, 2500.00, 'CREDIT',   1000002);
INSERT INTO TRANZACTII_C VALUES (1000004, SYSDATE - 11, 1800.00, 'DEBIT',    1000002);
INSERT INTO TRANZACTII_C VALUES (1000005, SYSDATE - 21,  400.00, 'CREDIT',   1000003);
INSERT INTO TRANZACTII_C VALUES (1000006, SYSDATE -  6,  250.00, 'DEBIT',    1000003);
INSERT INTO TRANZACTII_C VALUES (1000007, SYSDATE - 23, 4000.00, 'CREDIT',   1000004);
INSERT INTO TRANZACTII_C VALUES (1000008, SYSDATE -  7, 2500.00, 'DEBIT',    1000004);
INSERT INTO TRANZACTII_C VALUES (1000009, SYSDATE - 17,  700.00, 'CREDIT',   1000005);
INSERT INTO TRANZACTII_C VALUES (1000010, SYSDATE -  2,  500.00, 'TRANSFER', 1000005);
INSERT INTO TRANZACTII_C VALUES (1000011, SYSDATE - 27, 1600.00, 'CREDIT',   1000006);
INSERT INTO TRANZACTII_C VALUES (1000012, SYSDATE - 13, 1100.00, 'DEBIT',    1000006);
INSERT INTO TRANZACTII_C VALUES (1000013, SYSDATE - 16, 3000.00, 'CREDIT',   1000007);
INSERT INTO TRANZACTII_C VALUES (1000014, SYSDATE -  5, 1400.00, 'DEBIT',    1000007);
INSERT INTO TRANZACTII_C VALUES (1000015, SYSDATE - 19, 4500.00, 'CREDIT',   1000008);
INSERT INTO TRANZACTII_C VALUES (1000016, SYSDATE - 10, 2200.00, 'DEBIT',    1000008);
INSERT INTO TRANZACTII_C VALUES (1000017, SYSDATE - 25, 9000.00, 'CREDIT',   1000009);
INSERT INTO TRANZACTII_C VALUES (1000018, SYSDATE - 12, 3500.00, 'DEBIT',    1000009);
INSERT INTO TRANZACTII_C VALUES (1000019, SYSDATE - 15, 1100.00, 'CREDIT',   1000010);
INSERT INTO TRANZACTII_C VALUES (1000020, SYSDATE -  3,  800.00, 'TRANSFER', 1000010);

-- CARDURI_C
INSERT INTO CARDURI_C VALUES (1000001, '5425233430109903', DATE '2028-08-31', 'ACTIV',  1000001);
INSERT INTO CARDURI_C VALUES (1000002, '2222420000001113', DATE '2027-04-30', 'ACTIV',  1000002);
INSERT INTO CARDURI_C VALUES (1000003, '5425233430109929', DATE '2029-01-31', 'ACTIV',  1000003);
INSERT INTO CARDURI_C VALUES (1000004, '4917300800000000', DATE '2026-11-30', 'BLOCAT', 1000004);
INSERT INTO CARDURI_C VALUES (1000005, '4007702835532454', DATE '2028-07-31', 'ACTIV',  1000005);

-- ANGAJATI_C
INSERT INTO ANGAJATI_C VALUES (1000001, 'Lupsa Andrei',    'Director',           8500.00, 2);
INSERT INTO ANGAJATI_C VALUES (1000002, 'Chis Mirela',     'Consultant Bancar',  4500.00, 2);
INSERT INTO ANGAJATI_C VALUES (1000003, 'Bota Emil',       'Casier',             3200.00, 2);
INSERT INTO ANGAJATI_C VALUES (1000004, 'Suciu Maria',     'Analist Credite',    5200.00, 2);
INSERT INTO ANGAJATI_C VALUES (1000005, 'Ardelean Ionut',  'Responsabil Carduri',4800.00, 2);

-- ANGAJAT_CLIENT_C
INSERT INTO ANGAJAT_CLIENT_C VALUES (1000001, 1,  'DIRECTOR',  SYSDATE - 200);
INSERT INTO ANGAJAT_CLIENT_C VALUES (1000002, 3,  'GESTIONAR', SYSDATE - 130);
INSERT INTO ANGAJAT_CLIENT_C VALUES (1000002, 4,  'GESTIONAR', SYSDATE - 100);
INSERT INTO ANGAJAT_CLIENT_C VALUES (1000003, 5,  'GESTIONAR', SYSDATE -  70);
INSERT INTO ANGAJAT_CLIENT_C VALUES (1000004, 2,  'ANALIST',   SYSDATE - 160);

COMMIT;

ALTER SEQUENCE SEQ_CLIENT_C INCREMENT BY 10; SELECT SEQ_CLIENT_C.NEXTVAL FROM DUAL; ALTER SEQUENCE SEQ_CLIENT_C INCREMENT BY 1;
ALTER SEQUENCE SEQ_CONT_C   INCREMENT BY 10; SELECT SEQ_CONT_C.NEXTVAL   FROM DUAL; ALTER SEQUENCE SEQ_CONT_C   INCREMENT BY 1;
ALTER SEQUENCE SEQ_TRANZ_C  INCREMENT BY 20; SELECT SEQ_TRANZ_C.NEXTVAL  FROM DUAL; ALTER SEQUENCE SEQ_TRANZ_C  INCREMENT BY 1;

\end{lstlisting}
\screenshot{populate-cluj.png}{59 rânduri inserate pe nodul Cluj}

% ──────────────────────────────────────────────────────────────────────────────
\section{Furnizarea formelor de transparență}

\begin{lstlisting}
CREATE VIEW CLIENTI_GLOBAL AS
SELECT ci.ID_Client, ci.Nume, ci.Prenume, ci.CNP,
       cp.Email, cp.Telefon, cp.Scor_Credit
FROM   CLIENTI_ID@BUCHAREST_LINK ci
JOIN   CLIENTI_PROFIL cp ON ci.ID_Client = cp.ID_Client
UNION
SELECT ci.ID_Client, ci.Nume, ci.Prenume, ci.CNP,
       cp.Email, cp.Telefon, cp.Scor_Credit
FROM   CLIENTI_ID@CLUJ_LINK ci
JOIN   CLIENTI_PROFIL cp ON ci.ID_Client = cp.ID_Client
WHERE  ci.ID_Client NOT IN (
         SELECT ID_Client FROM CLIENTI_ID@BUCHAREST_LINK
       );

-- CONTURI_GLOBAL -- reconstructie fragmentare orizontala
CREATE VIEW CONTURI_GLOBAL AS
SELECT ID_Cont, IBAN, ID_Tip, Sold, Moneda, ID_Sucursala, ID_Client, 'B' AS Nod
FROM   CONTURI_B@BUCHAREST_LINK
UNION ALL
SELECT ID_Cont, IBAN, ID_Tip, Sold, Moneda, ID_Sucursala, ID_Client, 'C' AS Nod
FROM   CONTURI_C@CLUJ_LINK;

-- TRANZACTII_GLOBAL -- reconstructie fragmentare orizontala derivata
CREATE VIEW TRANZACTII_GLOBAL AS
SELECT ID_Tranzactie, Data, Suma, Tip_Tranzactie, ID_Cont_Sursa, 'B' AS Nod
FROM   TRANZACTII_B@BUCHAREST_LINK
UNION ALL
SELECT ID_Tranzactie, Data, Suma, Tip_Tranzactie, ID_Cont_Sursa, 'C' AS Nod
FROM   TRANZACTII_C@CLUJ_LINK;

-- trg_clienti_global
-- INSERT: generare ID global, inscrie in ambele noduri CLIENTI_ID si in CLIENTI_PROFIL
-- UPDATE: propaga nume/CNP catre ambele noduri, iar datele referitoare la profil pe Central
-- DELETE: sterge de pe ambele noduri si de pe central
CREATE OR REPLACE TRIGGER trg_clienti_global
INSTEAD OF INSERT OR UPDATE OR DELETE ON CLIENTI_GLOBAL
FOR EACH ROW
DECLARE
  v_id NUMBER(10);
BEGIN
  IF INSERTING THEN
    v_id := CASE WHEN :NEW.ID_Client IS NOT NULL THEN :NEW.ID_Client
                 ELSE SEQ_CLIENT_GLOBAL.NEXTVAL
            END;
    INSERT INTO CLIENTI_ID@BUCHAREST_LINK (ID_Client, Nume, Prenume, CNP)
      VALUES (v_id, :NEW.Nume, :NEW.Prenume, :NEW.CNP);
    INSERT INTO CLIENTI_ID@CLUJ_LINK (ID_Client, Nume, Prenume, CNP)
      VALUES (v_id, :NEW.Nume, :NEW.Prenume, :NEW.CNP);
    INSERT INTO CLIENTI_PROFIL (ID_Client, Email, Telefon, Scor_Credit)
      VALUES (v_id, :NEW.Email, :NEW.Telefon, NVL(:NEW.Scor_Credit, 500));

  ELSIF UPDATING THEN
    UPDATE CLIENTI_ID@BUCHAREST_LINK
      SET Nume=:NEW.Nume, Prenume=:NEW.Prenume, CNP=:NEW.CNP
      WHERE ID_Client = :OLD.ID_Client;
    UPDATE CLIENTI_ID@CLUJ_LINK
      SET Nume=:NEW.Nume, Prenume=:NEW.Prenume, CNP=:NEW.CNP
      WHERE ID_Client = :OLD.ID_Client;
    UPDATE CLIENTI_PROFIL
      SET Email=:NEW.Email, Telefon=:NEW.Telefon, Scor_Credit=:NEW.Scor_Credit
      WHERE ID_Client = :OLD.ID_Client;

  ELSIF DELETING THEN
    DELETE FROM CLIENTI_ID@BUCHAREST_LINK WHERE ID_Client = :OLD.ID_Client;
    DELETE FROM CLIENTI_ID@CLUJ_LINK     WHERE ID_Client = :OLD.ID_Client;
    DELETE FROM CLIENTI_PROFIL           WHERE ID_Client = :OLD.ID_Client;
  END IF;
END;
/

-- trg_conturi_global
-- INSERT:  verifica unicitate globala IBAN, ruteaza dupa ID_Sucursala
-- UPDATE:  update in-place pe acelasi fragment (Sold, IBAN, etc.)
-- DELETE:  delete in cascada pe acelasi fragment
CREATE OR REPLACE TRIGGER trg_conturi_global
INSTEAD OF INSERT OR UPDATE OR DELETE ON CONTURI_GLOBAL
FOR EACH ROW
DECLARE
  v_id  NUMBER(10);
  v_cnt NUMBER(1);
BEGIN
  IF INSERTING THEN
    -- Unicitate globala IBAN: verifica ambele fragmente
    SELECT COUNT(*) INTO v_cnt
    FROM (
      SELECT IBAN FROM CONTURI_B@BUCHAREST_LINK WHERE IBAN = :NEW.IBAN
      UNION ALL
      SELECT IBAN FROM CONTURI_C@CLUJ_LINK       WHERE IBAN = :NEW.IBAN
    ) WHERE ROWNUM = 1;
    IF v_cnt > 0 THEN
      RAISE_APPLICATION_ERROR(-20001,
        'IBAN ' || :NEW.IBAN || ' exista deja global (unicitate globala IBAN)');
    END IF;

    IF :NEW.ID_Sucursala = 1 THEN
      SELECT SEQ_CONT_B.NEXTVAL@BUCHAREST_LINK INTO v_id FROM DUAL;
      INSERT INTO CONTURI_B@BUCHAREST_LINK
        (ID_Cont, IBAN, ID_Tip, Sold, Moneda, ID_Sucursala, ID_Client)
        VALUES (v_id, :NEW.IBAN, :NEW.ID_Tip, NVL(:NEW.Sold, 0),
                NVL(:NEW.Moneda, 'RON'), 1, :NEW.ID_Client);
    ELSIF :NEW.ID_Sucursala = 2 THEN
      SELECT SEQ_CONT_C.NEXTVAL@CLUJ_LINK INTO v_id FROM DUAL;
      INSERT INTO CONTURI_C@CLUJ_LINK
        (ID_Cont, IBAN, ID_Tip, Sold, Moneda, ID_Sucursala, ID_Client)
        VALUES (v_id, :NEW.IBAN, :NEW.ID_Tip, NVL(:NEW.Sold, 0),
                NVL(:NEW.Moneda, 'RON'), 2, :NEW.ID_Client);
    END IF;

  ELSIF UPDATING THEN
    IF :OLD.ID_Sucursala = 1 THEN
      UPDATE CONTURI_B@BUCHAREST_LINK
        SET IBAN=:NEW.IBAN, ID_Tip=:NEW.ID_Tip, Sold=:NEW.Sold, Moneda=:NEW.Moneda
        WHERE ID_Cont = :OLD.ID_Cont;
    ELSE
      UPDATE CONTURI_C@CLUJ_LINK
        SET IBAN=:NEW.IBAN, ID_Tip=:NEW.ID_Tip, Sold=:NEW.Sold, Moneda=:NEW.Moneda
        WHERE ID_Cont = :OLD.ID_Cont;
    END IF;

  ELSIF DELETING THEN
    IF :OLD.ID_Sucursala = 1 THEN
      DELETE FROM TRANZACTII_B@BUCHAREST_LINK WHERE ID_Cont_Sursa = :OLD.ID_Cont;
      DELETE FROM CARDURI_B@BUCHAREST_LINK    WHERE ID_Cont       = :OLD.ID_Cont;
      DELETE FROM CONTURI_B@BUCHAREST_LINK    WHERE ID_Cont       = :OLD.ID_Cont;
    ELSE
      DELETE FROM TRANZACTII_C@CLUJ_LINK WHERE ID_Cont_Sursa = :OLD.ID_Cont;
      DELETE FROM CARDURI_C@CLUJ_LINK    WHERE ID_Cont       = :OLD.ID_Cont;
      DELETE FROM CONTURI_C@CLUJ_LINK    WHERE ID_Cont       = :OLD.ID_Cont;
    END IF;
  END IF;
END;
/

-- TRANZACTII_GLOBAL INSTEAD OF trigger
-- Ruteaza: ID_Cont_Sursa < 1000000 -> Bucharest; else Cluj
CREATE OR REPLACE TRIGGER trg_tranzactii_global
INSTEAD OF INSERT OR UPDATE OR DELETE ON TRANZACTII_GLOBAL
FOR EACH ROW
DECLARE
  v_id NUMBER(10);
BEGIN
  IF INSERTING THEN
    IF :NEW.ID_Cont_Sursa < 1000000 THEN
      SELECT SEQ_TRANZ_B.NEXTVAL@BUCHAREST_LINK INTO v_id FROM DUAL;
      INSERT INTO TRANZACTII_B@BUCHAREST_LINK
        (ID_Tranzactie, Data, Suma, Tip_Tranzactie, ID_Cont_Sursa)
        VALUES (v_id, NVL(:NEW.Data, SYSDATE), :NEW.Suma, :NEW.Tip_Tranzactie, :NEW.ID_Cont_Sursa);
    ELSE
      SELECT SEQ_TRANZ_C.NEXTVAL@CLUJ_LINK INTO v_id FROM DUAL;
      INSERT INTO TRANZACTII_C@CLUJ_LINK
        (ID_Tranzactie, Data, Suma, Tip_Tranzactie, ID_Cont_Sursa)
        VALUES (v_id, NVL(:NEW.Data, SYSDATE), :NEW.Suma, :NEW.Tip_Tranzactie, :NEW.ID_Cont_Sursa);
    END IF;

  ELSIF UPDATING THEN
    IF :OLD.ID_Cont_Sursa < 1000000 THEN
      UPDATE TRANZACTII_B@BUCHAREST_LINK
        SET Data=:NEW.Data, Suma=:NEW.Suma, Tip_Tranzactie=:NEW.Tip_Tranzactie
        WHERE ID_Tranzactie = :OLD.ID_Tranzactie;
    ELSE
      UPDATE TRANZACTII_C@CLUJ_LINK
        SET Data=:NEW.Data, Suma=:NEW.Suma, Tip_Tranzactie=:NEW.Tip_Tranzactie
        WHERE ID_Tranzactie = :OLD.ID_Tranzactie;
    END IF;

  ELSIF DELETING THEN
    IF :OLD.ID_Cont_Sursa < 1000000 THEN
      DELETE FROM TRANZACTII_B@BUCHAREST_LINK WHERE ID_Tranzactie = :OLD.ID_Tranzactie;
    ELSE
      DELETE FROM TRANZACTII_C@CLUJ_LINK WHERE ID_Tranzactie = :OLD.ID_Tranzactie;
    END IF;
  END IF;
END;
/

-- Verificare CLIENTI_GLOBAL
INSERT INTO CLIENTI_GLOBAL (Nume, Prenume, CNP, Email, Telefon, Scor_Credit)
VALUES ('Test', 'Global', '1990101400999', 'test.global@bank.ro', '0799000001', 650);
SELECT ID_Client, Nume, Prenume, Scor_Credit FROM CLIENTI_GLOBAL
  WHERE CNP = '1990101400999';

SELECT ID_Client FROM CLIENTI_ID@BUCHAREST_LINK WHERE CNP = '1990101400999';
SELECT ID_Client FROM CLIENTI_ID@CLUJ_LINK     WHERE CNP = '1990101400999';
ROLLBACK;

-- Verificare CONTURI_GLOBAL - transparenta orizontala
-- Utilizatorul interactioneaza doar cu CONTURI_GLOBAL, iar triggerul ruteaza catre
-- fragmentul corespunzator

INSERT INTO CONTURI_GLOBAL (IBAN, ID_Tip, Sold, Moneda, ID_Sucursala, ID_Client)
VALUES ('RO49GLOBTEST00000000001', 1, 1000, 'RON', 1, 1);
INSERT INTO CONTURI_GLOBAL (IBAN, ID_Tip, Sold, Moneda, ID_Sucursala, ID_Client)
VALUES ('RO49GLOBTEST00000000002', 1, 2500, 'RON', 2, 1);
COMMIT;

SELECT ID_Cont, IBAN, Sold, ID_Sucursala, Nod FROM CONTURI_GLOBAL
  WHERE IBAN IN ('RO49GLOBTEST00000000001', 'RO49GLOBTEST00000000002');

UPDATE CONTURI_GLOBAL SET Sold = 1500
  WHERE IBAN = 'RO49GLOBTEST00000000001';
COMMIT;

SELECT ID_Cont, IBAN, Sold, Nod FROM CONTURI_GLOBAL
  WHERE IBAN = 'RO49GLOBTEST00000000001';

-- Cleanup
DELETE FROM CONTURI_GLOBAL
  WHERE IBAN IN ('RO49GLOBTEST00000000001', 'RO49GLOBTEST00000000002');
COMMIT;
\end{lstlisting}

\subsection{Transparență verticală — CLIENTI\_GLOBAL}
\screenshot{trg-clienti-global.png}{CLIENTI\_GLOBAL returnează 10 clienți reconstituiți din ambele fragmente}
\screenshot{transparenta-verticala.png}{S1: INSERT via CLIENTI\_GLOBAL — rândul apare în CLIENTI\_ID@BUCHAREST\_LINK și @CLUJ\_LINK}

\subsection{Transparență orizontală — CONTURI\_GLOBAL și TRANZACTII\_GLOBAL}
\screenshot{trg-conturi-global.png}{CONTURI\_GLOBAL reconstituie fragmentele B și C — UNION ALL returnează 20 conturi cu coloana Nod}
\screenshot{transparenta-orizontala-exemplu-insert-conturi-global.png}{S2: INSERT via CONTURI\_GLOBAL — ID\_Sucursala=1 redirecționat în CONTURI\_B (Nod=B), ID\_Sucursala=2 în CONTURI\_C (Nod=C); UPDATE Sold propagat automat la fragmentul corect}

% ──────────────────────────────────────────────────────────────────────────────
\section{Asigurarea sincronizării datelor pentru relațiile replicate}

Tabela replicată este \texttt{TIPURI\_CONT} (4 înregistrări: Curent, Economii, Depozit, Credit),
stocată pe Central și propagată sincron la nodurile București și Cluj.
Sincronizarea se realizează în două etape: o populare inițială prin \texttt{MERGE}
și un trigger \texttt{AFTER INSERT OR UPDATE OR DELETE} pe tabela master care propagă
orice modificare în aceeași tranzacție spre ambele noduri replicate prin DB Links.

\begin{lstlisting}
MERGE INTO TIPURI_CONT@BUCHAREST_LINK dst
USING TIPURI_CONT src ON (dst.ID_Tip = src.ID_Tip)
WHEN MATCHED     THEN UPDATE SET dst.Denumire = src.Denumire, dst.Dobanda = src.Dobanda
WHEN NOT MATCHED THEN INSERT (ID_Tip, Denumire, Dobanda) VALUES (src.ID_Tip, src.Denumire, src.Dobanda);

MERGE INTO TIPURI_CONT@CLUJ_LINK dst
USING TIPURI_CONT src ON (dst.ID_Tip = src.ID_Tip)
WHEN MATCHED     THEN UPDATE SET dst.Denumire = src.Denumire, dst.Dobanda = src.Dobanda
WHEN NOT MATCHED THEN INSERT (ID_Tip, Denumire, Dobanda) VALUES (src.ID_Tip, src.Denumire, src.Dobanda);

COMMIT;

-- trg_replicate_tipuri_cont
CREATE OR REPLACE TRIGGER trg_replicate_tipuri_cont
AFTER INSERT OR UPDATE OR DELETE ON TIPURI_CONT
FOR EACH ROW
BEGIN
  IF INSERTING THEN
    INSERT INTO TIPURI_CONT@BUCHAREST_LINK (ID_Tip, Denumire, Dobanda)
      VALUES (:NEW.ID_Tip, :NEW.Denumire, :NEW.Dobanda);
    INSERT INTO TIPURI_CONT@CLUJ_LINK (ID_Tip, Denumire, Dobanda)
      VALUES (:NEW.ID_Tip, :NEW.Denumire, :NEW.Dobanda);

  ELSIF UPDATING THEN
    UPDATE TIPURI_CONT@BUCHAREST_LINK
      SET Denumire = :NEW.Denumire, Dobanda = :NEW.Dobanda
      WHERE ID_Tip = :OLD.ID_Tip;
    UPDATE TIPURI_CONT@CLUJ_LINK
      SET Denumire = :NEW.Denumire, Dobanda = :NEW.Dobanda
      WHERE ID_Tip = :OLD.ID_Tip;

  ELSIF DELETING THEN
    DELETE FROM TIPURI_CONT@BUCHAREST_LINK WHERE ID_Tip = :OLD.ID_Tip;
    DELETE FROM TIPURI_CONT@CLUJ_LINK     WHERE ID_Tip = :OLD.ID_Tip;
  END IF;
END;
/

-- Verificare -- inserarea unui nou tip de cont pe Central,

INSERT INTO TIPURI_CONT VALUES (5, 'Pensie', 2.50);
COMMIT;
SELECT ID_Tip, Denumire, Dobanda FROM TIPURI_CONT;
SELECT ID_Tip, Denumire, Dobanda FROM TIPURI_CONT@BUCHAREST_LINK;
SELECT ID_Tip, Denumire, Dobanda FROM TIPURI_CONT@CLUJ_LINK;
ROLLBACK;
\end{lstlisting}
\screenshot{replicare.png}{Sincronizare inițială: 4 rânduri MERGE în TIPURI\_CONT@BUCHAREST\_LINK și @CLUJ\_LINK}
\screenshot{replicare-trigger.png}{S6: INSERT pe TIPURI\_CONT master → apare automat în ambele replici (COUNT=1 pe fiecare nod)}

% ──────────────────────────────────────────────────────────────────────────────
\section{Asigurarea tuturor constrângerilor de integritate}

\subsection{Constrângeri de unicitate locală}

Constrângerile UNIQUE sunt definite direct în DDL-ul fiecărei tabele locale, garantând unicitatea
în cadrul fragmentului.

\begin{lstlisting}
-- Nod Bucuresti (BUCHAREST_USER)
CREATE TABLE TIPURI_CONT (
    ID_Tip    NUMBER(4)    PRIMARY KEY,
    Denumire  VARCHAR2(50) NOT NULL UNIQUE
    Dobanda   NUMBER(5,2)  NOT NULL
);
CREATE TABLE CLIENTI_ID (
    ID_Client  NUMBER(10)  PRIMARY KEY,
    Nume       VARCHAR2(100) NOT NULL,
    Prenume    VARCHAR2(100) NOT NULL,
    CNP        CHAR(13)    NOT NULL UNIQUE
);
CREATE TABLE CONTURI_B (
    ID_Cont  NUMBER(10)   PRIMARY KEY,
    IBAN     VARCHAR2(34) NOT NULL UNIQUE
    ...
);
CREATE TABLE CARDURI_B (
    ID_Card     NUMBER(10) PRIMARY KEY,
    Numar_Card  CHAR(16)   NOT NULL UNIQUE
);
-- Aceleasi constrangeri UNIQUE exista simetric pe nodul Cluj (CLUJ_USER)
-- pentru CONTURI_C, CARDURI_C, CLIENTI_ID, TIPURI_CONT
\end{lstlisting}

\subsection{Constrângeri de unicitate globală pe fragmente orizontale}

Unicitatea globală a \texttt{IBAN} între fragmentele \texttt{CONTURI\_B} și \texttt{CONTURI\_C}
este asigurată pe două niveluri: structural, prin ranguri de secvențe non-suprapuse, și
procedural, prin validare explicită în triggerul INSTEAD OF la momentul INSERT.

\begin{lstlisting}
-- Secvente cu domenii disjuncte garanteaza ID unic intre noduri
CREATE SEQUENCE SEQ_CONT_B  START WITH 1       INCREMENT BY 1 MAXVALUE 999999   NOCYCLE NOCACHE;
CREATE SEQUENCE SEQ_CONT_C  START WITH 1000001 INCREMENT BY 1 MAXVALUE 1999999  NOCYCLE NOCACHE;

-- Triggerul trg_iof_conturi_global verifica explicit ambele fragmente
-- inainte de INSERT, refuzand orice IBAN deja existent global
CREATE OR REPLACE TRIGGER trg_iof_conturi_global
INSTEAD OF INSERT OR UPDATE OR DELETE ON CONTURI_GLOBAL
FOR EACH ROW
DECLARE
  v_id  NUMBER(10);
  v_cnt NUMBER(1);
BEGIN
  IF INSERTING THEN
    SELECT COUNT(*) INTO v_cnt
    FROM (
      SELECT IBAN FROM CONTURI_B@BUCHAREST_LINK WHERE IBAN = :NEW.IBAN
      UNION ALL
      SELECT IBAN FROM CONTURI_C@CLUJ_LINK       WHERE IBAN = :NEW.IBAN
    ) WHERE ROWNUM = 1;
    IF v_cnt > 0 THEN
      RAISE_APPLICATION_ERROR(-20001,
        'IBAN ' || :NEW.IBAN || ' exista deja global (unicitate globala IBAN)');
    END IF;
    -- continua cu INSERT pe fragmentul corespunzator ID_Sucursala ...
  END IF;
END;

-- Test: INSERT cu IBAN duplicat intre fragmente:
INSERT INTO CONTURI_GLOBAL (IBAN, ID_Tip, Sold, Moneda, ID_Sucursala, ID_Client)
VALUES ('RO49GLOB0000001000000001', 1, 0, 'RON', 1, 1);
-- ORA-20001: IBAN RO49GLOB0000001000000001 exista deja global (unicitate globala IBAN)
ROLLBACK;
\end{lstlisting}

\screenshot{unicitate-globala.png}{Test constrangeri de unicitate globala}

\subsection{Constrângeri de unicitate globală pe fragmente verticale}

Unicitatea globală a \texttt{CNP} pe fragmentele verticale ale \texttt{CLIENTI} este asigurată
de triggerul INSTEAD OF pe \texttt{CLIENTI\_GLOBAL}: inserarea în \texttt{CLIENTI\_ID@BUCHAREST\_LINK}
eșuează dacă CNP-ul există deja pe orice nod (constrângere UNIQUE locală).

\begin{lstlisting}
-- Triggerul trg_iof_clienti_global insereaza simultan
-- pe ambele noduri. Daca CNP-ul exista deja pe oricare nod, INSERT esueaza:
CREATE OR REPLACE TRIGGER trg_iof_clienti_global
INSTEAD OF INSERT OR UPDATE OR DELETE ON CLIENTI_GLOBAL
FOR EACH ROW
BEGIN
  IF INSERTING THEN
    INSERT INTO CLIENTI_ID@BUCHAREST_LINK (ID_Client, Nume, Prenume, CNP)
      VALUES (v_id, :NEW.Nume, :NEW.Prenume, :NEW.CNP);
    INSERT INTO CLIENTI_ID@CLUJ_LINK (ID_Client, Nume, Prenume, CNP)
      VALUES (v_id, :NEW.Nume, :NEW.Prenume, :NEW.CNP);
  END IF;
END;

-- Test: INSERT cu CNP duplicat :
INSERT INTO CLIENTI_GLOBAL (Nume, Prenume, CNP, Email, Telefon, Scor_Credit)
VALUES ('Test', 'Dup', '1900101400001', 'dup@bank.ro', '0700000000', 500);
-- ORA-00001: unique constraint (BUCHAREST_USER.SYS_C...) violated on CNP
ROLLBACK;
\end{lstlisting}

\screenshot{constrangere-unicitate-globala-pe-fragmente-verticale.png}{Constrangere de unicitate globala pe fragmente verticale}


\subsection{Constrângeri de cheie primară}

Cheile primare sunt definite local în DDL-ul fiecărei tabele. Unicitatea globală între noduri
este garantată prin secvențe cu domenii disjuncte

\begin{lstlisting}
CREATE TABLE CONTURI_B    (ID_Cont       NUMBER(10) PRIMARY KEY, ...);
CREATE TABLE TRANZACTII_B (ID_Tranzactie NUMBER(10) PRIMARY KEY, ...);
CREATE TABLE CARDURI_B    (ID_Card       NUMBER(10) PRIMARY KEY, ...);
CREATE TABLE ANGAJATI_B   (ID_Angajat    NUMBER(10) PRIMARY KEY, ...);
CREATE TABLE CLIENTI_ID   (ID_Client     NUMBER(10) PRIMARY KEY, ...);

CREATE TABLE SUCURSALE    (ID_Sucursala  NUMBER(4)  PRIMARY KEY, ...);
CREATE TABLE TIPURI_CONT  (ID_Tip        NUMBER(4)  PRIMARY KEY, ...);
CREATE TABLE CLIENTI_PROFIL (ID_Client   NUMBER(10) PRIMARY KEY, ...);
CREATE TABLE CREDITE      (ID_Credit     NUMBER(10) PRIMARY KEY, ...);
CREATE TABLE PLATI_RATE   (ID_Plata      NUMBER(10) PRIMARY KEY, ...);

-- Unicitate globala garantata prin ranguri disjuncte:
-- Bucuresti: SEQ_CONT_B  [1       .. 999999]
-- Cluj:      SEQ_CONT_C  [1000001 .. 1999999]
-- Central:   SEQ_CLIENT_GLOBAL [2000001 .. MAX]
\end{lstlisting}

\subsection{Constrângeri de cheie externă}

Cheile externe locale sunt implementate ca \texttt{FOREIGN KEY} DDL. Cheile externe cross-nod
(Central $\rightarrow$ noduri locale) sunt implementate prin triggere \texttt{BEFORE INSERT}.

\begin{lstlisting}
-- FK locale pe nodul Bucuresti
CREATE TABLE CONTURI_B (
    ...
    CONSTRAINT fk_conturi_b_cli FOREIGN KEY (ID_Client) REFERENCES CLIENTI_ID(ID_Client),
    CONSTRAINT fk_conturi_b_tip FOREIGN KEY (ID_Tip)    REFERENCES TIPURI_CONT(ID_Tip)
);
CREATE TABLE TRANZACTII_B (
    ...
    CONSTRAINT fk_tranz_b_cont FOREIGN KEY (ID_Cont_Sursa) REFERENCES CONTURI_B(ID_Cont)
);
CREATE TABLE CARDURI_B (
    ...
    CONSTRAINT fk_card_b_cont FOREIGN KEY (ID_Cont) REFERENCES CONTURI_B(ID_Cont)
);
CREATE TABLE ANGAJAT_CLIENT_B (
    ...
    CONSTRAINT fk_ac_b_ang FOREIGN KEY (ID_Angajat) REFERENCES ANGAJATI_B(ID_Angajat),
    CONSTRAINT fk_ac_b_cli FOREIGN KEY (ID_Client)  REFERENCES CLIENTI_ID(ID_Client)
);
CREATE TABLE PLATI_RATE (
    ...
    CONSTRAINT fk_plata_credit FOREIGN KEY (ID_Credit) REFERENCES CREDITE(ID_Credit)
);
\end{lstlisting}

\begin{lstlisting}
-- FK cross-nod: CREDITE (Central) -> CLIENTI_ID (Bucharest/Cluj)
-- Nu se poate implementa ca DDL FOREIGN KEY intre instante Oracle diferite;
-- se implementeaza ca trigger BEFORE INSERT care verifica nodul remote.
CREATE OR REPLACE TRIGGER trg_bef_ins_credite
BEFORE INSERT ON CREDITE
FOR EACH ROW
DECLARE
  v_count NUMBER;
BEGIN
  SELECT COUNT(*) INTO v_count
  FROM   CLIENTI_ID@BUCHAREST_LINK
  WHERE  ID_Client = :NEW.ID_Client;
  IF v_count = 0 THEN
    RAISE_APPLICATION_ERROR(-20001,
      'ID_Client ' || :NEW.ID_Client || ' nu exista in CLIENTI_ID');
  END IF;
END;

-- Test: INSERT credit pentru client inexistent:
INSERT INTO CREDITE (ID_Credit, Suma_Totala, Rata_Lunara, ID_Client)
VALUES (9999, 10000, 500, 99999);
-- ORA-20001: ID_Client 99999 nu exista in CLIENTI_ID
ROLLBACK;
\end{lstlisting}

\screenshot{constrangere-de-cheie-externa.png}{Constrangere de cheie externa cross-nod}

\subsection{Constrângeri de validare locală (CHECK)}

\begin{lstlisting}
CREATE TABLE CONTURI_B (
    ...
    CONSTRAINT chk_conturi_b_suc CHECK (ID_Sucursala = 1),        -- puritate fragment
    CONSTRAINT chk_moneda_b      CHECK (Moneda IN ('RON','EUR','USD')),
    CONSTRAINT chk_sold_b        CHECK (Sold >= 0)
);
CREATE TABLE TRANZACTII_B (
    ...
    CONSTRAINT chk_tranz_b_suma CHECK (Suma > 0),
    CONSTRAINT chk_tranz_b_tip  CHECK (Tip_Tranzactie IN ('DEBIT','CREDIT','TRANSFER'))
);
CREATE TABLE CARDURI_B (
    ...
    CONSTRAINT chk_card_b_status CHECK (Status IN ('ACTIV','BLOCAT','EXPIRAT'))
);
CREATE TABLE ANGAJATI_B (
    ...
    CONSTRAINT chk_ang_b_suc CHECK (ID_Sucursala = 1),
    CONSTRAINT chk_ang_b_sal CHECK (Salariu > 0)
);
-- Cluj: constrangeri simetrice cu CHECK (ID_Sucursala = 2)

-- Central: validare domeniu
CREATE TABLE TIPURI_CONT (
    ...
    CONSTRAINT chk_dobanda    CHECK (Dobanda >= 0 AND Dobanda <= 100)
);
CREATE TABLE CLIENTI_PROFIL (
    ...
    CONSTRAINT chk_scor  CHECK (Scor_Credit BETWEEN 0 AND 999),
    CONSTRAINT chk_email CHECK (Email LIKE '%@%.%')
);
CREATE TABLE CREDITE (
    ...
    CONSTRAINT chk_credit_suma CHECK (Suma_Totala > 0),
    CONSTRAINT chk_credit_rata CHECK (Rata_Lunara > 0)
);
CREATE TABLE SUCURSALE (
    ...
    CONSTRAINT chk_sucursala_oras CHECK (Oras IN ('Bucuresti', 'Cluj'))
);
\end{lstlisting}

\subsection{Constrângeri de validare globală pe fragmente diferite}

Constrângerile de validare care implică atribute distribuite pe noduri diferite sunt implementate
prin triggere \texttt{BEFORE INSERT OR UPDATE} care accesează nodul remote prin DB Link.

\begin{lstlisting}
-- INSERT direct local vizibil in vederea globala
INSERT INTO CONTURI_B@BUCHAREST_LINK
  (ID_Cont, IBAN, ID_Tip, Sold, Moneda, ID_Sucursala, ID_Client)
VALUES (SEQ_CONT_B.NEXTVAL@BUCHAREST_LINK,
        'RO49GLOBVERIF00000001', 1, 500.00, 'RON', 1, 1);
COMMIT;
SELECT ID_Cont, IBAN, Nod FROM CONTURI_GLOBAL
  WHERE IBAN = 'RO49GLOBVERIF00000001';

DELETE FROM CONTURI_B@BUCHAREST_LINK WHERE IBAN = 'RO49GLOBVERIF00000001';
COMMIT;

INSERT INTO CLIENTI_GLOBAL (Nume, Prenume, CNP, Email, Telefon, Scor_Credit)
VALUES ('Verif', 'Test', '5000101410001', 'verif.test@bank.ro', '0700000001', 700);
COMMIT;
SELECT ID_Client, Nume FROM CLIENTI_ID@BUCHAREST_LINK WHERE CNP = '5000101410001';
SELECT ID_Client, Nume FROM CLIENTI_ID@CLUJ_LINK     WHERE CNP = '5000101410001';
DELETE FROM CLIENTI_GLOBAL WHERE CNP = '5000101410001';
COMMIT;
\end{lstlisting}

\screenshot{local-visible-global.png}{S4: INSERT direct în CONTURI\_B → vizibil imediat în CONTURI\_GLOBAL cu Nod=B}
\screenshot{global-to-both-locals.png}{S5: INSERT via CLIENTI\_GLOBAL → rândul apare în ambele CLIENTI\_ID locale}

% ══════════════════════════════════════════════════════════════════════════════
\setcounter{section}{0}
\part{Modul Implementare Aplicație (N3)}
% ══════════════════════════════════════════════════════════════════════════════

\section{Modul gestiune date la nivel local}

Ruta \texttt{/local/<branch>} (Bucharest / Cluj) permite:
\begin{itemize}
  \item Vizualizarea conturilor și tranzacțiilor locale
  \item Adăugarea unui cont nou (\texttt{INSERT INTO CONTURI\_B/C} direct pe nodul local)
  \item Ștergerea unui cont
  \item Adăugarea unei tranzacții
\end{itemize}

\screenshot{bucharest-view.png}{Pagina /local/bucharest — lista conturi și tranzacții din CONTURI\_B}
\screenshot{add-cont-form.png}{Formular adăugare cont nou la sucursala București}
\screenshot{cont-added.png}{Confirmare cont adăugat — rândul apare în tabel}
\screenshot{cluj-view.png}{Pagina /local/cluj — lista conturi și tranzacții din CONTURI\_C}

\section{Modul vizualizare date la nivel global}

Ruta \texttt{/global} se conectează la \texttt{GLOBAL\_USER} și afișează:
\begin{itemize}
  \item \texttt{CLIENTI\_GLOBAL} — toți clienții reconstituiți din fragmente verticale
  \item \texttt{CONTURI\_GLOBAL} — toate conturile cu coloana Nod (B/C)
  \item Totaluri per sucursală (nr. conturi, sold total)
\end{itemize}

\screenshot{global-view.png}{Pagina /global — CLIENTI\_GLOBAL (10 clienți) și CONTURI\_GLOBAL (20 conturi, B+C)}

\section{Local → Global: vizualizare modificări locale la nivel global}

\subsection{Fragment orizontal (CONTURI)}
Ruta \texttt{/demo/local-to-global}: insert direct în \texttt{CONTURI\_B} sau \texttt{CONTURI\_C},
afișare before/after în \texttt{CONTURI\_GLOBAL}.

\screenshot{local-to-global.png}{Formular: selectare sucursală, introducere IBAN, sold}
\screenshot{before-after-horizontal.png}{Before: 10 conturi B; After: 11 conturi B}
\screenshot{before-after-horizontal-global.png}{Rândul nou apare în CONTURI\_GLOBAL}

\subsection{Fragment vertical (CLIENTI)}
Insert direct în \texttt{CLIENTI\_ID} pe nodul local → verificare în \texttt{CLIENTI\_GLOBAL}.

\screenshot{vertical-before-after.png}{Insert local în CLIENTI\_ID → rândul apare în CLIENTI\_ID@CLUJ_LINK}
\screenshot{vertical-before-after-global.png}{Insert local în CLIENTI\_ID → rândul apare în CLIENTI\_GLOBAL}

\subsection{Relații replicate (TIPURI\_CONT)}
Insert pe master \texttt{TIPURI\_CONT} → replicat automat pe replica → vizibil în ambele noduri.

\screenshot{replication-demo.png}{Insert TIPURI\_CONT master → COUNT=1 pe BUCHAREST și CLUJ replica}

\section{Global → Local: verificare propagare operații globale la nivel local}

\subsection{Fragment vertical (CLIENTI)}
Ruta \texttt{/demo/global-to-local}: INSERT via \texttt{CLIENTI\_GLOBAL} →
triggerul INSTEAD OF propagă la ambele \texttt{CLIENTI\_ID} locale.

\screenshot{global-to-local-propagation.png}{Clientul apare în CLIENTI\_ID@BUCHAREST\_LINK și CLIENTI\_ID@CLUJ\_LINK}

\subsection{Fragment orizontal (CONTURI)}
INSERT via \texttt{CONTURI\_GLOBAL} cu \texttt{ID\_Sucursala=1} → rândul apare în \texttt{CONTURI\_B};
cu \texttt{ID\_Sucursala=2} → rândul apare în \texttt{CONTURI\_C}.

\screenshot{horizontal-routing.png}{INSERT via CONTURI\_GLOBAL cu Sucursala=1 → rândul apare în CONTURI\_B (Nod=B)}

\end{document}
